\documentclass{../template/note}

\author{Oliver West}
\title{Groups}
\date{}

\begin{document}
    
    \maketitle

    \tableofcontents

    \lecture{Lecture 1 -- 10/10/25}

    \section{Examples and definitions}

    One can think about groups in two ways - \emph{algebra} and \emph{symmetry}.

    \begin{figure}[ht]
        \centering
        \incfig{0.3\textwidth}{equilateral-triangle-symmetry}
        \caption{Symmetries of the equilateral triangle.}
        \label{fig:equilateral-triangle-symmetry}
    \end{figure}
    
    \noindent We also include the \emph{identity} symmetry of doing nothing.

    \begin{figure}[ht]
        \centering
        \incfig{0.4\textwidth}{equilateral-triangle-all-symmetries}
        \caption{All 6 symmetries of an equilateral triangle.}
        \label{fig:equilateral-triangle-all-symmetries}
    \end{figure}

    \begin{remark*}
        Can we generalise? How many symmetries does a regular $n$-gon have?
    \end{remark*}
    
    Things get more interesting if we \emph{compose} -- \emph{multiply?} our symmetries.

    \begin{figure}[ht]
        \centering
        \incfig{0.4\textwidth}{equilateral-triangle-composition-example}
        \caption{Composing our symmetries gives another symmetry!}
        \label{fig:equilateral-triangle-composition-example*}
    \end{figure}
    
    Let's tabulate our properties.
    \begin{itemize}
        \item There is an \emph{identity} symmetry.
        \item Every symmetry has an \emph{inverse}.
        \begin{itemize}
            \item The inverse of the identity is itself.
            \item The inverse of a clockwise rotation is the corresponding counter-clockwise rotation.
            \item The inverse of a reflection is itself!
        \end{itemize}
        \item It is possible to compose symmetries at all!
        \item \emph{Associativity}: Composition of symmetries is associative -- it doesn't matter how we bracket when we compose our symmetries.
        \item We \emph{don't} necessarily have \emph{commutativity}.
        
        \begin{figure}[ht]
            \centering
            \incfig{0.4\textwidth}{equilateral-triangle-commutativity-failure}
            \caption{When we perform the operations from \cref{fig:equilateral-triangle-composition-example*} in the opposite order, the outcome is different.}
            \label{fig:equilateral-triangle-commutativity-failure}
        \end{figure}
        
    \end{itemize}

    Time to go back to algebra.

    \begin{definition}
        A \emph{binary operation} on a set $X$ is a function
        \begin{align*}
            \cdot\; \colon X \times X &\to X \\
            (x, y) &\mapsto x \cdot y
        \end{align*}
    \end{definition}
    
    \begin{definition}
        A \emph{group} is a triple $(G, \cdot, e)$ where
        \begin{itemize}
            \item $G$ is a set;
            \item $\cdot$ is a binary operation;
            \item and $e$ is an element of $G$.
        \end{itemize}
        such that the following \emph{group axioms} hold:
        \begin{itemize}
            \item \emph{Closure}. For all $a$, $b \in G$, $a \cdot b \in G$ (we actually already know this because it's part of the definition of a binary operation).
            \item \emph{Associativity}. For all $a$, $b$, $c \in G$, $(a \cdot b) \cdot c = a \cdot (b \cdot c)$.
            \item \emph{Identity}. For all $a \in G$, $a \cdot e = a$.
            \item \emph{Right inverses}. For all $a \in G$, there exists $b \in G$ such that $a \cdot b = e$. 
        \end{itemize}        
    \end{definition}
    
    \begin{example*}
        The symmetries of an equilateral triangle form a group.
    \end{example*}
    
    \begin{exercise}
        Show that $(\ZZ, +, 0)$ is a group.
    \end{exercise}
    
    \begin{proposition}
        Let $(G, \cdot, e)$ be a group, and let $a, b, b', e'$ be elements of $G$.
        \begin{enumerate}
            \item If $a \cdot b = e$, then $b \cdot a = e$ -- right inverses are left inverses;
            \item $e \cdot a = a$ -- the identity is also a left identity;
            \item If $a$, $b$ and $b'$ are such that $a \cdot b = a \cdot b' = e$, then $b = b'$ -- inverses are unique;
            \item If $a \cdot e' = a$ then $e' = e$ -- the identity is unique.
        \end{enumerate}
        
        
    \end{proposition}

    \lecture{Lecture 2 -- 13/10/25}

    \begin{proposition}
        Let $(G, \cdot, e)$ be a group, and let $a, b, b', e'$ be elements of $G$.
        \begin{enumerate}
            \item If $a \cdot b = e$, then $b \cdot a = e$ -- right inverses are left inverses;
            \item $e \cdot a = a$ -- the identity is also a left identity;
            \item If $a$, $b$ and $b'$ are such that $a \cdot b = a \cdot b' = e$, then $b = b'$ -- inverses are unique;
            \item If $a \cdot e' = a$ then $e' = e$ -- the identity is unique.
        \end{enumerate}
        
        
    \end{proposition}

    \begin{proof}
        We simply prove each part sequentially.
        \begin{enumerate}
            \item Suppose $a \cdot b = e$. Then
            \begin{align*}
                b &= b \cdot e && \text{by identity} \\
                &= b \cdot (a \cdot b) && \text{by assumption} \\
                &= (b \cdot a) \cdot b && \text{by associativity} \\
                \implies b \cdot b^{-1} &= ((b \cdot a ) \cdot b) \cdot b^{-1} && \text{where $b^{-1}$ is an inverse of $b$} \\
                &= (b \cdot a) \cdot (b \cdot b^{-1}) && \text{by associativity} \\
                \implies e &= (b \cdot a) \cdot e && \text{by definition of the inverse} \\
                &= b \cdot a && \text{by definition of the identity}
            \end{align*}
            \item Let $a^{-1}$ be an inverse of $a$, then by 1, $a \cdot a^{-1} = e = a^{-1} \cdot a$. Now
            \begin{align*}
                e \cdot a &= (a \cdot a^{-1}) \cdot a \\
                &= a \cdot (a^{-1} \cdot a) \\
                &= a \cdot e \\
                &= a
            \end{align*}
            \item Suppose $a \cdot b = a \cdot b' = e$. Then
            \begin{align*}
                b' &= e \cdot b' && \text{by 2} \\
                &= (b \cdot a) \cdot b' && \text{by 1} \\
                &= b \cdot (a \cdot b') \\
                &= b \cdot e \\
                &= b
            \end{align*}
            \item Let $a^{-1}$ be the inverse of $a$, so $a^{-1} \cdot a = e$.
            \begin{align*}
                a &= a \cdot e' \\
                \implies a^{-1} \cdot a &= a^{-1} \cdot (a \cdot e') \\
                &= (a^{-1} \cdot a) \cdot e' \\
                &= e \cdot e' \\
                &= e'
            \end{align*}
            Hence $e = e'$. Note we only needed $e'$ to be a `second inverse' for one element!
        \end{enumerate}
    \end{proof}
    
    A corollary of the uniqueness of inverses along with right inverses also being left inverses is that $a$ is the inverse of $a^{-1}$.

    \begin{definition}[Powers of elements]
        Fix $a \in G$, then
        \begin{itemize}
            \item $a^0 = e$;
            \item For $n \in \NN$, $a^n = a^{n-1} \cdot a$.
            \item For $n < 0$, $a^n = (a^{-1})^{-n}$.
        \end{itemize}        
    \end{definition}
    
    \begin{proposition}
        For $a \in G$, $n, m \in \ZZ$,
        \begin{itemize}
            \item $a^m \cdot a^n = a^{m+n}$;
            \item $(a^m)^n = a^{mn}$.
        \end{itemize}
    \end{proposition}
    
    \begin{proposition}
        For $a$, $b \in G$, $(a\cdot b)^{-1} = b^{-1}\cdot a^{-1}$.
    \end{proposition}
    
    \begin{definition}[Abelian groups]
        For a group $(G, \cdot, e)$, if $a \cdot b = b \cdot a$ for all $a$, $b \in G$, then $G$ is \emph{abelian}.
    \end{definition}
    
    \begin{example*}[The trivial group]
        The \emph{trivial group} has $G = \{e\}$ and $e \cdot e = e$
    \end{example*}
    
    

    \begin{example*}[Additive groups]
        $(\ZZ, +, 0)$, $(\QQ, +, 0)$, $(\RR, +, 0)$, $(\CC, +, 0)$ -- all abelian. $(\NN, +, 0)$ is \emph{not} a group.
    \end{example*}

    \begin{example*}[Multiplicative groups]
        $(\QQ, \times, 1)$ is not a group -- but $(\QQ \setminus \{0\}, \times, 1)$ is (abelian). As are $(\RR \setminus \{0\}, \times, 1)$ and $(\CC \setminus \{0\}, \times, 1)$.
    \end{example*}
    
    \begin{definition}
        The \emph{order} of a group $(G, \cdot, e)$ is the number of elements of $G$, $|G|$. If $|G| < \infty$, then $(G, \cdot, e)$ is \emph{finite}.
    \end{definition}
    
    \begin{example*}[Finite groups]
        \begin{itemize}
            \item For $n \in \NN$, let $C_n = \{z: z \in \CC, z^n = 1\}$. Then $(C_n, \times, 1)$ is an abelian group.
            \item For $n \in \NN$, let $Z_n = \{0, \dots, n-1\}$. Let $a +_n b$ be the remainder when $a + b$ is divided by $n$. Then $(Z_n, +_n, 0)$ is an abelian group.
        \end{itemize}
        
        
    \end{example*}

    \lecture{Lecture 3 -- 15/10/25}

    \subsection{Symmetric groups}

    \begin{definition}
        Let $X$, $Y$ be sets. A \emph{bijection} is a map $f: X \to Y$ that has an inverse $g: Y \to X$ such that $f \circ g = \id_Y$ and $g \circ f = \id_X$.
    \end{definition}
    
    \begin{definition}
        A bijection from a set $X$ to itself is called a \emph{permutation}.
    \end{definition}
    
    \begin{definition}
        Given a set $X$, the set $\Sym(X)$ is the set of permutations of the set.
    \end{definition}
    
    \begin{lemma}[Associativity of function composition]
        Consider maps of sets:

        \[\begin{tikzcd}
            W & X & Y & Z
            \arrow["f", from=1-1, to=1-2]
            \arrow["g", from=1-2, to=1-3]
            \arrow["h", from=1-3, to=1-4]
        \end{tikzcd}\]
        Then $(h \circ g) \circ f = h \circ (g \circ f)$.
        \label{associativity-of-function-composition}
    \end{lemma}
    
    \begin{proof}
        For any $w \in W$, $((h \circ g) \circ f)(w) = h(g(f(w))) = (h \circ (g \circ f))(w)$.
    \end{proof}
    
    And that's all we need to show to make our next group!

    \begin{proposition}
        For any set $X$, $(\Sym(X), \circ, \id_X)$ is a group.
    \end{proposition}
    
    \begin{proof}
        We must verify the four group axioms.
        \begin{itemize}
            \item \emph{Closure}. $(f \circ g)^{-1} = g^{-1} \circ f^{-1}$, so the composition of two permutations is a permutation.
            \item \emph{Associativity}. By \cref{associativity-of-function-composition}.
            \item \emph{Identity}. $\id_X$.
            \item \emph{Inverse}. By definition: we are working with bijections.
        \end{itemize}
        
        
    \end{proof}
    
    \begin{definition}
        If $X = \{1, \dots, n\}$, we write $S_n$ for $(\Sym(X), \circ, \id_X)$.
    \end{definition}
    
    \begin{example*}
        \begin{itemize}
            \item $S_3$, for example, is the group of ways to arrange 3 flowerpots on the windowsill.
            \item $S_{52}$ is the group of ways to shuffle a deck of cards.
        \end{itemize}
    \end{example*}
    
    \noindent Fairly clear that $|S_n| = n!$.

    \begin{remark*}[Notation]
        Writing $(G, \cdot, e)$ is cumbersome. Let's just write $G$ and leave the rest as implied. If we really want to emphasise that $e$ is the identity, we can always write $e_G$, and $\cdot_G$ for $\cdot$. Finally we might be so apathetic that we even write $ab$ instead of $a \cdot b$.
    \end{remark*}
    
    \subsection{Subgroups}

    We might want to focus on smaller groups within a group -- for example, $\ZZ$ inside $\RR$, or the rotations of a triangle instead of all its symmetries.

    \begin{definition}
        Let $G$ be a group, and $H \subseteq G$. If $H$ satisfies
        \begin{itemize}
            \item $e \in H$;
            \item for any $a$, $b \in H$, $ab \in H$;
            \item for any $a \in H$, $a^{-1} \in H$
        \end{itemize}
        then we say $H$ is a \emph{subgroup} of $G$. We write $H \leq G$.
        
    \end{definition}
    
    \begin{remark*}
        If $G$ is a group and $H$ is a subgroup, then $H$ is itself a group.
    \end{remark*}
    
    \begin{example*}
        \begin{itemize}
            \item Every group $G$ has itself and $\{e\}$ as subgroups. Any other subgroup is called \emph{proper}.
            \item $\ZZ \leq \QQ \leq \RR \leq \CC$.
            \item For $n \in \NN_0$, let $n\ZZ = \{nk : k \in \ZZ\}$ -- the set of multiples of $n$. Then $n\ZZ \leq \ZZ$. In fact these are the only subgroups of $\ZZ$
            
        \end{itemize}
        
        
    \end{example*}
    
    \begin{proposition}
        If $H \leq \ZZ$, $H = n\ZZ$ for some $n \in \NN_0$.
    \end{proposition}
    
    \begin{proof}
        If $H = \{0\}$, then $H = 0\ZZ$. If $H \neq \{0\}$, then let $n$ be the smallest positive integer in $H$. By induction, $nk \in H$ for all $k \in \ZZ$, so $n\ZZ \leq H$. Now suppose for contradiction there exists $h \in H$ such that $h \notin n\ZZ$. Dividing $h$ by $n$ gives $h = nq + r$ with $0 \leq r < n$, and in particular by closure $r \in H$. Furthermore, $r \neq 0$ since $h \notin n\ZZ$. But this contradicts the minimality of $n$. Hence $H = n\ZZ$ as required.
    \end{proof}
    
    \lecture{Lecture 4 -- 17/10/25}

    \begin{proposition}
        If $H$, $K \leq G$, then $H \cap K \leq G$
    \end{proposition}
    
    \begin{proof}
        (Example sheet 1)
    \end{proof}
    
    \begin{remark*}
        We can generalise to any collection of subgroups.
    \end{remark*}
    
    \begin{definition}
        Let $X$ be a subset of a group $G$. Then
        \[
            \gen{X} := \bigcap_{X \leq H \leq G} H
        \]
        is the \emph{subgroup generated by $X$} -- the intersection of all subgroups of $G$ that contain $X$. Intuition: the smallest subgroup of $G$ containing $X$.
    \end{definition}
    
    \begin{definition}
        If $\gen X = G$, we say that $X$ \emph{generates} $G$ or is a \emph{generating set} for $G$. If $X$ generates $G$, it means that every $g \in G$ can be written as
        \[
            x_1^{\pm 1} \cdot x_2^{\pm 1} \cdots x_n^{\pm 1}
        \]
        for some $n \geq 0$, where every $x_i \in X$.
    \end{definition}
    
    \subsection{Geometric examples}

    Let $\CC$ be the plane, equipped with the usual notion of distance. 

    DIAGRAM: complex plane, w = w1 + iw2, z = z1 + iz2, line between them, labelled |z-w| = an expression due to Pythagoras

    \begin{definition}
        For any subset $X \subseteq \CC$, an \emph{isometry} of $X$ is a bijection $f$ from $X$ to itself that preserves distance -- for any $x$, $y \in X$, $|f(x)-f(y)|=|x-y|$.
    \end{definition}
    
    \begin{proposition}[Isometry groups in $\CC$]
        Let $X \subseteq \CC$. The set of isometries of $X$, $\Isom(X)$, is a subgroup of $\Sym(X)$. In particular, $\Isom(X)$ is a group.
    \end{proposition}
    
    \begin{proof}
        Let $f$, $g \in \Isom(X)$, and $x$, $y \in X$.
        \begin{itemize}
            \item \emph{Identity}. Clearly $\id_X$ is an isometry, so it is in $\Isom(X)$ as required.
            \item \emph{Closure}. Clear that the composition of isometries is an isometry by the transitivity of equality.
            \item \emph{Inverse}. If $f$ is an isometry,
            \begin{align*}
                |f^{-1}(x) - f^{-1}(y)| &= |f \circ f^{-1}(x) - f \circ f^{-1}(y)| \\
                &= |x - y|
            \end{align*}
        \end{itemize}
    \end{proof}
    
    \begin{definition}
        Let $n \geq 3$ and let $X_n \subseteq \CC$ be the $n$-gon with vertices $\{\exp(\frac{2 \pi i k}{n}) : k = 0, \dots, n-1\}$.

        DIAGRAM: $6$ points equally spaced around the unit circle in the complex plane, one point at $n=1$, label `e.g. $n=6$', hexagon drawn.

        Then the \emph{$n^\text{th}$ dihedral group}, $D_{2n} := \Isom(X_n)$. So for example, $D_6$ is the symmetry group of an equilateral triangle.
    \end{definition}
    
    \begin{theorem}
        The size of the $n^\text{th}$ dihedral group is $2n$.
    \end{theorem}
    
    First, we need some geometric lemmas.

    \begin{lemma}[Kite lemma]
        Let $x_1$, $x_2$, $y_1$, $y_2 \in \CC$. If $|y_1-x_1| = |y_2-x_1$ and $|y_1-x_2| = |y_2-x_2|$, then $x_2-x_1$ is perpendicular to $y_2-y_1$.
    \end{lemma}
    
    \begin{proof}
        DIAGRAM: Kite: diagonals y1y2 and x1x2. Whimsical kite tail attached to x2. Equal lengths of x1y1 and x1y2 and x2y1 and x2y2 indicated. Diagonals drawn. Point $z$ labelled at intersection.

        By symmetry, $\angle X_1 Z Y_1 = \angle X_1 Z Y_2$. But $Y_1Y_2$ is a straight line, so $\angle X_1 Z Y_1 + \angle X_1 Z Y_2 = \pi$. Hence $\angle X_1 Z Y_1 = \angle X_1 Z Y_2 = \frac \pi 2$ as required.
    \end{proof}
    
    \begin{lemma}[Three point lemma]
        Let $X \subseteq \CC$ and $f \in \Isom(X)$. Suppose $x_1$, $x_2$, $x_3 \in X$ distinct and not collinear, such that $f(x_i) = x_i$ for $i=1,2,3$. Then $f = \id_X$.
    \end{lemma}
    
    \begin{proof}
        Suppose $f(y) \neq y$ for some $y \in X$. Then
        \begin{align*}
            |f(y) - x_i| &= |f(y) - f(x_i)| \\
            &= |y - x_i| && \text{since } f \in \Isom(X)
        \end{align*}
        Now apply the kite lemma with $y_1 = y$, $y_2 = f(y)$ and letting $x_1$, $x_2$ range over the $x_i$ to get $f(y) - y$ perpendicular to $x_i - x_j$ for any $i \neq j$. Then, since $f(y)-y \neq 0$,  $x_i - x_j$ are all parallel, so $x_1$, $x_2$, $x_3$ are collinear.

        DIAGRAM: ????
        
    \end{proof}
    
    \begin{remark*}
        There is equally an $(n+1)$-point lemma valid in $\RR^n$.
    \end{remark*}
    
    \lecture{Lecture 5 -- 20/10/25}

    \begin{proof}[Proof of theorem]
        We first define two particular elements of $D_{2n}$:
        \begin{align*}
            r: z &\mapsto e^{\frac{2 \pi i}{n}}z && \text{rotation} \\
            s: z &\mapsto \overline{z} && \text{reflection} \\
        \end{align*}

        \begin{figure}[ht]
            \centering
            \incfig{0.3\textwidth}{dihedral-group-operations}
            \caption{Graphical depictions of $r$ and $s$}
            \label{fig:dihedral-group-operations}
        \end{figure}      

        We will now further claim that $D_{2n} = \{e, r, \dots, r^{n-1}, s, rs, \dots, r^{n-1}s\}$. In particular, $\{r, s\}$ generates $D_{2n}$.

        First, it is necessary to show that $r$, $s \in D_{2n}$. For some $x$, $y \in \CC$,
        \begin{itemize}
            \item $|r(x) - r(y)| = |e^{\frac{2 \pi i}{n}} x - e^{\frac{2 \pi i}{n}} y| = |e^{\frac{2 \pi i}{n}}||x - y| = |x-y|$, so $r$ is an isometry. Furthermore, $r(e^{\frac{2 \pi i k}{n}}) = e^{\frac{2 \pi i (k+1)}{n}}$, so $r$ sends vertices to other vertices, so $r$ really is a permutation of $X_n$.
            \item $|s(x) - s(y)| = |\overline{x} - \overline{y}| = |\overline{x-y}| = |x-y|$, so $s$ is an isometry. To see that $s$ preserves the vertices of $X_n$, we can either proceed algebraically or appeal to the symmetry of $X_n$ in the real axis.
        \end{itemize}
        
        From this it follows that $\{e, r, \dots, r^{n-1}, s, rs, \dots, r^{n-1}s\} \subseteq D_{2n}$ by closure. It remains to show that these are all the elements. Let $f$ arbitrary in $D_{2n}$. Let $x=1$, $y=e^{\frac{2\pi i}{n}}$, $z=e^{-\frac{2\pi i}{n}}$ (DIAGRAM?). Since $f \in D_{2n}$, $f(x)$ is some vertex of $X_n$ -- that is, $f(x) = e^{\frac{2 \pi i k}{n}}$ for some $k \in \{0, \dots, n-1\}$. Then $r^{-k} \circ f(x) = 1$. Since $r^{-k} \circ f$ is an isometry, we can also see $r^{-k} \circ f(y)$ and $r^{-k} \circ f(z)$ are themselves $y$ and $z$ in some order.

        \begin{enumerate}[label=\protect\circled{\arabic*}]
            \item $r^{-k} \circ f(y) = y$ and $r^{-k} \circ f(z) = z$. Then, by the 3 point lemma, $r^{-k} \circ f = e$. So $f = r^k \in \{e, r, \dots, r^{n-1}, s, rs, \dots, r^{n-1}s\}$.
            \item $r^{-k} \circ f(y) = z$ and $r^{-k} \circ f(z) = y$. Then $s \circ r^{-k} \circ f(y) = y$ and $s \circ r^{-k} \circ f(z) = z$, so $s \circ r^{-k} \circ f = e$. Hence $f = r^ks$, since $s$ is self-inverse, so $f \in \{e, r, \dots, r^{n-1}, s, rs, \dots, r^{n-1}s\}$, as required.
        \end{enumerate}
        
        Thus $D_{2n} \subseteq \{e, r, \dots, r^{n-1}, s, rs, \dots, r^{n-1}s\}$. Finally, it is necessary to show that these elements are all distinct. One can simply verify that $x$ and $y$ are taken to different points by each of these. This concludes the proof.
        
    \end{proof}
    
    To understand the group operation on $D_{2n}$, it is desirable to understand all the different ways to multiply $r$ and $s$.

    \begin{lemma}[Dihedral relation]
        For $r$, $s \in D_{2n}$ as above,
        \[
            sr = r^{-1}s
        \]
    \end{lemma}
    
    \begin{proof}
        One can appeal to intuition or just check for a small number of cases by the 3 point lemma.
    \end{proof}
    
    \lecture{Lecture 6 -- 22/10/25}
    
    \subsection{Homomorphisms}

    Some groups are not set-theoretically equal, but nevertheless have the same structure -- for example, $\Sym(\text{deck of cards})$ and $\Sym(\text{52 undergraduate mathematicians})$.

    \begin{definition}[Homomorphism]
        A map between groups $\phi: G \to H$ is called a \emph{homomorphism} if, for any $g$, $g' \in G$,
        \[
            \phi(g \cdot g') = \phi(g) \cdot \phi(g')
        \]
    \end{definition}
    
    \begin{example*}
        \begin{enumerate}[label=\protect\circled{\arabic*}]
            \item For any groups $G$, $H$, the map $\phi: G \to H$, $g \mapsto e_H$ is the \emph{trivial homomorphism}.
            \item If $H \leq G$, the map $i: H \to G$, $h \mapsto h$ is the \emph{inclusion homomorphism}
            \item If $nk = m$, the map $\phi: C_m \to C_n$, $g \mapsto g^k$ is a homomorphism.
            \item Since $\det(AB) = \det(A) \det(B)$, $\det: \GL_2 \RR \to (\RR \setminus \{0\}, \times)$, $A \mapsto \det A$ is a homomorphism.
        \end{enumerate}
    \end{example*}
    
    \begin{lemma}
        If $\phi: G \to H$ is a homomorphism, then
        \begin{enumerate}
            \item $\phi(e_G) = \phi(e_H)$, and
            \item $\phi(g^{-1}) = \phi(g)^{-1}$ for all $g \in G$.
        \end{enumerate}
        \label{homomorphism-identity-inverse}   
    \end{lemma}
    
    \begin{proof}
        \begin{enumerate}
            \item $\phi(e_G) \cdot \phi(e_G) = \phi(e_G \cdot e_G) = \phi(e_G)$, so $\phi(e_G) = e_H$, by uniqueness of identities.
            \item $\phi(g) \cdot \phi(g^{-1}) = \phi(g \cdot g^{-1}) = \phi(e_G) = e_H$, so $\phi(g)^{-1} = \phi(g^{-1})$, by uniqueness of inverses.
        \end{enumerate}        
    \end{proof}
    
    \begin{definition}[Isomorphisms]
        If a homomorphism $\phi: G \to H$ is also a bijection, then $\phi$ is called an \emph{isomorphism}, and we write $G \cong H$.
    \end{definition}
    
    \begin{example*}
        \begin{enumerate}
            \item $\ZZ_n \cong C_n$ -- we have the isomorphism $\phi: \ZZ_n \to C_n$, $k \mapsto e^{2\pi i k /n}$.
            \item $(\RR, +, 0) \cong (\RR_{> 0}, \times, 1)$ -- we have the isomorphism $\phi: (\RR, +, 0) \to (\RR_{> 0}, \times, 1)$, $x \mapsto \exp(x)$.
        \end{enumerate}
    \end{example*}
    
    \begin{lemma}
        \begin{enumerate}
            \item If $\phi: G \to H$ is an isomorphism, so is $\phi^{-1}$.
            \item If $\phi: G \to H$ and $\psi: H \to K$ are homomorphisms, so is $\psi \circ \phi$.
            \item $\cong$ is an equivalence relation.
        \end{enumerate}
    \end{lemma}
    
    \begin{proof}
        Exercise.
    \end{proof}
    
    We saw earlier that every subgroup induces an inclusion homomorphism. Furthermore, every homomorphism induces a subgroup.

    \begin{definition}[Image and kernel]
        Let $\phi: G \to H$ be a homomorphism.
        \begin{itemize}
            \item The \emph{image} of $\phi$, $\im \phi $, is $\{h \in H: h = \phi(g) \text{ for some } g \in G\}$. (DIAGRAM)
            \item The \emph{kernel} of $\phi$, $\ker \phi$, is $\{g \in G: \phi(g) = e_H\}$.
        \end{itemize}        
    \end{definition}
    
    \begin{proposition}
        If $\phi: G \to H$ is a homomorphism, then
        \begin{enumerate}
            \item $\im \phi \leq H$;
            \item $\ker \phi \leq G$.
        \end{enumerate}
    \end{proposition}
    
    \begin{proof}
        \begin{enumerate}
            \item $e_H = \phi(e_G) \in \im \phi$, so we have the identity. $\phi(g_1)\phi(g_2) = \phi(g_1g_2) \in \im \phi$, so we have closure. Finally, by \cref{homomorphism-identity-inverse}, $\phi(g)^{-1} = \phi(g^{-1}) \in \im \phi$, so we have inverse.
            \item $\phi(e_G) = e_H$, so $e_G \in \ker \phi$, so we have the identity. For $g_1, g_2 \in \ker \phi$, $\phi(g_1g_2) = \phi(g_1)\phi(g_2) = e_H \cdot e_H = e_H$, so $g_1g_2 \in \ker \phi$, so we have closure. Finally, by \cref{homomorphism-identity-inverse}, $\phi(g^{-1}) = \phi(g)^{-1} = e_H^{-1} = e_H$, so we have inverse.
        \end{enumerate}
    \end{proof}
    
    \begin{proposition}
        \label{surj-inv-crits}
        Let $\phi: G \to H$ be a homomorphism.
        \begin{enumerate}
            \item $\phi$ is \emph{surjective} if and only if $\im \phi = H$;
            \item $\phi$ is \emph{injective} if and only if $\ker \phi = \{e_G\}$.
        \end{enumerate}
    \end{proposition}
    
    \lecture{Lecture 7 -- 24/10/25}

    \begin{proof}
        The surjectivity statement is very clear. For the injectivity, we will show each direction as follows:
        \begin{itemize}
            \item[$\implies$] If $\phi$ is injective, then there is only one element that maps to $e_H$, i.e. the kernel must have one element. That element clearly must be $e_G$.
            \item[$\impliedby$] Suppose $\ker \phi = \{e_G\}$, and $\phi(g_1) = \phi(g_2)$. Then $\phi(g_1g_2^{-1}) = \phi(g_1)\phi(g_2)^{-1} = e_H$. Hence, $g_1g_2^{-1} \in \ker \phi \implies g_1g_2^{-1} = e_G \implies g_1 = g_2$, which was what we wanted.
        \end{itemize}
    \end{proof}
    
    Combining the two parts of the proposition, we see that a homomorphism $\phi$ is an isomorphism iff and only if $\im \phi = H$ and $\ker \phi = \{e_G\}$.

    \subsection{Cyclic groups}

    \begin{definition}[Cyclic groups]
        A group $G$ is \emph{cyclic} if there exists $g \in G$ such that $G = \{g^k: k \in \ZZ\} = \gen g$. $g$ is called a \emph{generator} of $G$.
    \end{definition}
    
    \begin{example*}
        Let us have some examples.
        \begin{itemize}
            \item $C_n = \{z \in \CC: z^n = 1\}$ is cyclic, with generator $e^{\frac{2\pi i}n}$;
            \item $\ZZ$ is cyclic, with generator 1;
            \item $\ZZ_n$ is cyclic, since $\ZZ_n \cong C_n$.
        \end{itemize}
    \end{example*}
    
    \begin{theorem}
        If $G$ is cyclic, then either
        \begin{itemize}
            \item $G \cong C_n$ for some $n$;
            \item or $G \cong \ZZ$.
        \end{itemize}
    \end{theorem}
    
    \begin{proof}
        Let $G$ be a cyclic group with generator $g$. Let
        \begin{align*}
            S = \{k \in \NN: g^k = e\}
        \end{align*}
        and let
        \begin{align*}
            n = \begin{cases}
                \min S & S \neq \emptyset \\
                \infty & \text{o/w}
            \end{cases}
        \end{align*}
        We now split into cases.

        \emph{Case 1.} $n = \infty$. Define
        \begin{align*}
            \phi: \ZZ &\to G \\
            k &\mapsto g^k
        \end{align*}
        which we would like to be an isomorphism. Since $\phi(k+l) = g^{k+l} = g^kg^l = \phi(k)\phi(l)$, $\phi$ is a homomorphism. Since $G$ is cyclic, every element of $G$ has the form $g^k$ for some $k$, so $\phi$ is surjective. To prove injectivity, suppose for contradiction $0 \neq k \in \ker \phi$ -- since $\ker \phi \leq \ZZ$, we may replace $k$ by $-k$ if necessary, and assume $k > 0$. But then $k \in S$, which is a contradiction. The result follows.
        
        \emph{Case 2.} $n < \infty$. Define
        \begin{align*}
            \phi: \ZZ_n &\to G \\
            k &\mapsto g^k
        \end{align*}
        as before. Since $k+l = qn + (k +_n l)$ for some $q \in \ZZ$,
        \begin{align*}
            \phi(k)\phi(l) &= g^kg^l = g^{k+l} = g^{qn+(k+_nl)} \\
            &= g^{k+_nl} = \phi(k +_n l)
        \end{align*}
        since $g^n = e$. Hence $\phi$ is a homomorphism. Since $G$ is cyclic, every element is $g^k$ for some $k \in \ZZ$. Dividing $k$ by $n$, we obtain
        \begin{align*}
            k = qn + r
        \end{align*}
        for $q \in \RR$ and $r \in \ZZ_n$. Hence,
        \begin{align*}
            g^k = g^{qn+r} = g^r = \phi(r)
        \end{align*}
        as desired. To prove injectivity, suppose that $\phi(k) = e$ for some $k \in \ZZ_n$, that is $0 \leq k < n$. On the other hand, $k \in S$, so, if $k \neq 0$, $k \geq n$ by minimality of $n$, a contradiction. Hence $k=0$, so $\ker \phi = \{0\}$ and $\phi$ is injective as required. Hence $G \cong \ZZ_n \cong C_n$ as required.
    \end{proof}
    
    Because of the theorem, we will write $C_n$ for any cyclic group of order $n$, and write $C_\infty \cong \ZZ$.

    \begin{definition}[Order of an element]
        For any group $G$ and $g \in G$, let
        \begin{align*}
            \gen g = \{g^k : k \in \ZZ\} \leq G
        \end{align*}
        Since $\gen g$ is cyclic, $\gen g \cong C_n$ for some $n \in \NN \cup \{\infty\}$. This is the \emph{order} of $g$, $o(g)$ or $|g|$.        
    \end{definition}
    
    \subsection{Dihedral groups revisited}

    We would like to be able to tell if a group is dihedral just by checking some algebraic facts.

    \begin{lemma}
        \label{gen-dih-relation}
        Suppose $x$, $y \in G$ satisfy the dihedral relation $xy = yx^{-1}$. Then
        \begin{align*}
            yx^l = x^{-l}y
        \end{align*}
        for any $l \in \ZZ$.        
    \end{lemma}
    
    \begin{proof}
        First, take $l>0$, then
        \begin{align*}
            yx^{-l} = x^ly
        \end{align*}
        by induction. Furthermore, if $y^2 = e$, then
        \begin{align*}
            yx^l = y (x^l y) y = y(yx^{-l})y = x^{-l}y
        \end{align*}
        so that $yx^l = x^{-l}y$ for any $l \in \ZZ$.
    \end{proof}
    
    \begin{lemma}
        Let $G$ be a group, and $a$, $b \in G$ such that
        \begin{itemize}
            \item $a^n = e$ for some $n \geq 3$;
            \item $b^2 = e$;
            \item $ab = ba^{-1}$.
        \end{itemize}
        Then $\phi(r) = a$ and $\phi(s) = b$ defines a homomorphism $\phi: D_{2n} \to G$ sending $r$ to $a$ and $s$ to $b$.
    \end{lemma}
    
    \begin{proof}
        There are four cases to check.
        \begin{enumerate}[label=\protect\circled{\arabic*}]
            \item $\phi(r^k)\phi(r^l) = a^ka^l = a^{k+l} = a^{k+_nl} = \phi(r^{k+_nl}) = \phi(r^{k+l})$;
            \item $\phi(r^k)\phi(r^ls) = a^ka^lb = a^{k+l}b = a^{k+_nl}b = \phi(r^{k+_nl}s) = \phi(r^{k+l}s)$;
            \item $\phi(r^ks)\phi(r^l) = a^kba^l = a^{k-l}b = \phi(r^{k-l}s) = \phi(r^ksr^l)$ by \cref{gen-dih-relation}.
            \item $\phi(r^ks)\phi(r^ls) = a^k b a^l b = a^{k-l} b^2 = a^{k-l} = \phi(r^{k-l}) = \phi(r^{k-l}s^2) = \phi(r^k s r^l s)$, also by \cref{gen-dih-relation}.
        \end{enumerate}
    \end{proof}
    
    \begin{proposition}
        Suppose $G$ has a generating set $\{a, b\}$ such that
        \begin{itemize}
            \item $a^n = e$ for some $n \geq 3$;
            \item $b^2 = e$;
            \item $ab = ba^{-1}$;
            \item $|G| = 2n$.
        \end{itemize}
        Then $G \cong D_{2n}$.
    \end{proposition}
    
    \begin{proof}
        By the lemma, there is a homomorphism $\phi: D_{2n} \to G$. $\phi$ is surjective since $\{a, b\}$ generate $G$. Since $|D_{2n}| = 2n = |G|$, it must be a bijective. The result follows.
    \end{proof}
    
    
    
    


    \lecture{Lecture 8 -- 27/10/25}

    \section{Lagrange's theorem}

    This very important theorem helps us think about the orders of groups and their subgroups.

    \begin{theorem}[Lagrange]
        If $H \leq G$ and $|G| < \infty$, $|H| \mid |G|$.
    \end{theorem}
    
    \subsection{Cosets}

    \begin{definition}{Left cosets}
        Let $H \leq G$ and $g \in G$. The corresponding \emph{left coset} is
        \[
            gH = \{gh : h \in H\} \subseteq G
        \]
        The set of all left cosets of $H$ is denoted by $G/H$.
    \end{definition}
    
    \begin{remark*}
        We may equally define right cosets, $Hg$ -- the set of right cosets of $H$ in $G$ is denoted by $H \backslash G$.
    \end{remark*}
    
    \begin{lemma}
        If $H \leq G$, the left cosets of $H$ partition $G$ -- that is, every element of $G$ is in exactly one coset. In other words, the following hold:
        \begin{enumerate}
            \item
            \begin{align*}
                G = \bigcup_{gH \in G/H} gH
            \end{align*}
            \item If $g_1H \cap g_2H \neq \emptyset$, then $g_1H = g_2H$. 
        \end{enumerate} 
    \end{lemma}
    
    \begin{proof}
        \begin{enumerate}
            \item For any $g \in G$, $g = g \cdot e \in gH$, and so clearly $\bigcup gH = G$.
            \item Suppose $g_1H \cap g_2H \neq \emptyset$, so there exists $k$ with $k \in g_1H$, $k \in g_2H$, so $k = g_1h_1 = g_2h_2$ for some $h_1$, $h_2 \in H$. Thus $g_1 = g_2 (h_2h_1^{-1}) \in g_2H$.
            
            Furthermore, for any $H \in H$, $g_1h = g_2((h_2h_1^{-1})h) \in g_2H$, so $g_1H \subseteq g_2H$. But the same reasoning would show $g_2H \subseteq g_1H$, and so we are done.
        \end{enumerate}
    \end{proof}
    
    DIAGRAM: G as a rectangle made of rectangles labelled with cosets.

    But turns out things are even nicer than this -- all the cosets have the same size.

    \begin{lemma}
        Let $H \leq G$ -- then there is a bijection $H \to gH$ for any $g \in G$.
    \end{lemma}
    
    \begin{proof}
        We claim the map $H \to gH$, $h \mapsto gh$ is a bijection. Observe we have an inverse $gh \to H$, $gh \mapsto g^{-1}(gh) = h$. This concludes the proof.
    \end{proof}
    
    \begin{corollary}
        All left cosets have the same size.
    \end{corollary}

    DIAGRAM: G as a rectangle made of equally sized rectangles.

    \begin{definition}[Index of a subgroup]
        The \emph{index} of a subgroup $H \leq G$ is $|G : H| := |G / H|$, the number of cosets of $H$ in $G$.
    \end{definition}
    
    \begin{remark*}
        One could show that all of these results continue to hold for right cosets.
    \end{remark*}
    
    \begin{theorem}[Lagrange]
        If $|G| < \infty$, then
        \[
            |G| = |G:H||H|
        \]
        for any $H \leq G$.
    \end{theorem}
    
    \begin{proof}
        Since left cosets partition $G$, $|G|$ is the sum of the sizes of all the cosets in $G$, which is $|G : H||H|$.
    \end{proof}
    
    \subsection{Consequences of Lagrange's theorem}

    \begin{corollary}
        If $|G| < \infty$ and $g \in G$, then $o(g) \mid |G|$.
    \end{corollary}
    
    \begin{proof}
        $o(g) = |\gen{g}|$.
    \end{proof}
    
    \begin{corollary}
        If $|G| < \infty$ and $g \in G$, then $g^{|G|} = e$.
        \label{element-to-order-is-identity}
    \end{corollary}
    
    \begin{proof}
        Since $o(g) \mid |G|$, $g^{|G|}$ = $(g^{o(g)})^{G/o(g)} = e^{G/o(g)} = e$.
    \end{proof}
    
    \begin{corollary}
        If $|G|$ is prime, then $G$ is cyclic.
        \label{prime-groups-cyclic}
    \end{corollary}
    
    \begin{proof}
        Let $g \in G$, $g \neq e$. Then $o(g) | |G|$, so, since $|G|$ prime, $o(g) = |G|$ (since $o(g) \neq e$). Therefore, $\gen g = G$.
    \end{proof}
    
    In fact, Lagrange's theorem implies an important result from number theory -- the Fermat-Euler theorem.

    \begin{definition}
        The \emph{Euler totient function} is given by
        \[
            \varphi(n) = |\{x \in \ZZ_n : \gcd(x,n)=1\}|
        \]
    \end{definition}
    
    \lecture{Lecture 9 -- 29/10/25}

    Let $\times_n$ denote multiplication modulo $n$ (on $\ZZ_n$). Recall that $x \in \ZZ_n$ has a multiplicative inverse modulo $n$ if and only if $\gcd(x, n) = 1$. Thus,
    \[
        \ZZ_n^\times = \{x \in \ZZ_n : \gcd(x, n) = 1\}
    \]
    is an abelian group under multiplication.

    \begin{theorem}[Fermat-Euler]
        Let $x, n \in \ZZ$ with $\gcd(x, n) = 1$. Then
        \[
            x^{\varphi(n)} \equiv 1 \pmod n
        \]
    \end{theorem}
    
    \begin{proof}
        By \cref{element-to-order-is-identity}, $x^{|\ZZ_n^\times|} = 1$ for any $x \in \ZZ_n^\times$. But $|\ZZ_n^\times| = \varphi(n)$, so $x^{\varphi(n)} = 1$ in $\ZZ_n^\times$, and so $x^{\varphi(n)} \equiv 1 \pmod n$ as required.
    \end{proof}
    
    \section{Group actions}

    Groups become groups of symmetries when we make them act on symmetric objects.

    \begin{definition}
        An \emph{action} of a group $G$ on a set $X$ is a map
        \begin{align*}
            G \times X &\to X \\
            (g, x) &\mapsto gx
        \end{align*}
        such that
        \begin{itemize}
            \item $ex = x$ for all $x \in X$;
            \item $(gh)x = g(hx)$ for all $g, h \in G$, $x \in X$.
        \end{itemize}        
    \end{definition}

    \begin{example*}
        \begin{itemize}
            \item For any group $G$ and a set $X$, $gx = x$ for all $g \in G$ defines the \emph{trivial action}.
            \item $\Sym(X) \acts X$ for any set $X$ by $fx = f(x)$ (remember $f$ is a permutation).
            \item If $G \acts X$ and $H \leq G$, then $H \acts X$, so, for example, $\Isom(\CC) \acts \CC$.
            \item $D_{2n}$ acts on $X_n$ (the regular $n$-gon), whether `filled in' or just the edges or vertices.
            \item Every group $G$ acts on itself by the \emph{left regular action} -- just normal group multiplication on the left.           
        \end{itemize}
        
        
    \end{example*}
    
    \begin{theorem}
        An action of a group $G$ on a set $X$ is equivalent to a homomorphism
        \[
            \phi: G \to \Sym(X)
        \]
        \label{actions-are-homomorphisms}
    \end{theorem}
    
    \begin{proof}
        Suppose $G \acts X$, and consider
        \begin{align*}
            t_g: X &\to X \\
            x &\mapsto gx
        \end{align*}
        for any $g \in G$. Then $t_{g^{-1}} \circ t_g(x) = t_{g^{-1}}(gx) = g^{-1}(gx) = ex = x$, and so $t_g$ has an inverse. Thus $t_g$ is a bijection and in particular a permutation, so $t_g \in \Sym(X)$. Let us then define
        \begin{align*}
            \phi: G &\to \Sym(X) \\
            g &\mapsto t_g
        \end{align*}
        We now wish to show that $\phi$ is a homomorphism. Observe
        \begin{align*}
            t_g \circ t_h (x) &= t_g(hx) \\
            &= ghx \\
            &= t_{gh}(x)
        \end{align*}
        This concludes the proof that every action corresponds to a homomorphism. Let us now prove the other direction. Given a homomorphism
        \[
            \phi: G \to \Sym(X)
        \]   
        we define an action $G \acts X$ by $gx = \phi(g)(x)$. Let us now check that this is in fact an action:
        \begin{itemize}
            \item $ex = \phi(e)x = \id_X(x) = x$;
            \item $(gh)x = \phi(gh)(x) = \phi(g) \circ \phi(h)(x) = g(hx)$.
        \end{itemize}
        Thus we are done.           
    \end{proof}
    
    \begin{theorem}[Cayley]
        Every group $G$ is isomorphic to a subgroup of $\Sym(X)$ for some $X$. Furthermore, if $|G| < \infty$, we may require $|X| < \infty$.
    \end{theorem}
    
    \begin{proof}
        Let $X = G$, and consider the left regular action $G \acts X=G$. By \cref{actions-are-homomorphisms}, this is equivalent to a homomorphism
        \[
            \phi: G \to \Sym(X)
        \]
        Let $H = \im \phi \leq \Sym(X)$, so we may view $\phi$ as a homomorphism $G \to H$, and we claim this is in fact an isomorphism. It is surjective by the definition of $H$, so it remains to show it is injective. For this we will claim $\ker \phi = \{e\}$. Observe $g \in \ker \phi \implies \phi(g) = \id_X$, so $g\gamma = \gamma$ for all $\gamma \in G$. From this it is clear that $g = e$. Thus $G \cong H \leq \Sym(X)$, which was what we wanted. Since $G = X$, it is certainly true that $|G| < \infty \implies |X| < \infty$.
    \end{proof}
    
    It turns out actions are extremely useful and we can use them to learn a great deal about groups.

    \begin{definition}[Orbits and stabilisers]
        Let $G \acts X$ and $x \in X$.
        \begin{itemize}
            \item The \emph{orbit} of $x$ is the set $Gx = \{gx: g \in G\}$.
            \item The \emph{stabiliser} of $x$ is the set $\Stab_G(x) = G_x = \{g \in G: gx = x\}$.
        \end{itemize}
        If $Gx = X$ for any $x$, we say the action is \emph{transitive}. If every element $g \neq e \in G$ has some $x \in X$ such that $gx \neq x$, then $G \acts X$ is called \emph{faithful}.
    \end{definition}
    
    \begin{remark*}
        $G \acts X$ is faithful if the associated homomorphism is injective $\iff$ the kernel is trivial.
    \end{remark*}

    \lecture{Lecture 10 -- 31/10/25}
    
    \begin{proposition}
        Suppose $G \acts X$, then
        \begin{itemize}
            \item $\Stab_G(x) \leq G$ for any $x \in X$;
            \item The orbits $\{Gy: y \in X\}$ form a partition of $X$.
        \end{itemize}
        \label{stab-subgroup-orbits-partition}
    \end{proposition}
    
    \begin{proof}
        \begin{itemize}
            \item We simply check that the stabiliser satisfies the definition of a subgroup.
            \begin{itemize}
                \item \emph{Closure under multiplication}. If $g$, $h \in \Stab_G(x)$, $(gh)x$ = $g(hx) = gx = x$, so $gh \in \Stab_G(x)$;
                \item \emph{Closure under inversion}. If $g \in \Stab_G(x)$, $g^{-1}x = g^{-1}(gx) = (g^{-1}g)x = ex = x$, so $g^{-1} \in \Stab_G(x)$;
                \item \emph{Identity}. $ex = x$ by definition, so $e \in \Stab_G(x)$.
            \end{itemize}
            Hence $\Stab_G(x) \leq G$ as required.
            \item We must show every element is in an orbit, and if two orbits overlap, they are the same orbit.
            \begin{itemize}
                \item Every element is in its own orbit -- this shows the first part.
                \item Now suppose $Gx_1 \cap Gx_2 \neq \emptyset$, so we have $y \in Gx_1$ and $\in Gx_2$, which means $y = g_1x_1 = g_2x_2$ for some $g_1$, $g_2 \in G$. Then
                \begin{align*}
                    x_1 &= (g_1^{-1}g_1)x_1 \\
                    &= g_1^{-1} (g_1x_1) \\
                    &= g_1^{-1} (g_2x_2) \\
                    &= (g_1^{-1}g_2) x_2 \in Gx_2
                \end{align*}
                In fact, for any $gx_1 \in Gx_1$, $gx_1 = g(g_1^{-1}g_2)x_2 \in Gx_2$, so $Gx_1 \subseteq Gx_2$, and likewise $Gx_2 \subseteq Gx_1$, which concludes the proof.
            \end{itemize}
        \end{itemize}
    \end{proof}
    
    \begin{example*}
        Consider $D_{2n} \acts X_n$, the regular $n$-gon. Let us attempt to find the orbits and stabilisers associated with this action in the case $x=1$

        Clearly $r^js(x) = r^j(x) = e^{2 \pi i j/n}$, and so we can take $x$ to any other element -- that is, $D_{2n}x = X_n$. From this and \cref{stab-subgroup-orbits-partition}, it's clear that $D_{2n}x = X_n$ for any $x \in X_n$.  A similar calculation makes clear that $\Stab_G(x) = \{e, s\}$.
    \end{example*}
    
    \begin{theorem}[Orbit-stabiliser theorem]
        Suppose $G \acts X$, and let $x \in X$. The formula $g \Stab_G(x) \mapsto gx$ gives a well-defined bijection $G/\Stab_G(x) \to Gx$ -- \emph{cosets of the stabiliser are in bijection with orbits}.
    \end{theorem}
    
    \begin{corollary}
        If $G \acts X$ and $x \in X$, then
        \[
            |G| = |Gx||\Stab_G(x)|
        \]
        \label{orbit-stabiliser-theorem}
    \end{corollary}
    
    \begin{proof}[Proof of corollary]
        The theorem gives that $|Gx|$ = $|G/\Stab_G(x)|$, since they are in bijection. But $|G/\Stab_G(x)| = |G|/|\Stab_G(x)|$ by Lagrange, which concludes the proof.
    \end{proof}
    
    \begin{example*}
        Consider $D_{2n} \acts X_n$. We saw that, if $x \in X_n$, then $|D_{2n}| = n$ and $|\Stab_{D_{2n}}(x)| = 2$, so $|D_{2n}| = 2n$. This proof is circular, but this circularity could be resolved if we had approached things in a different way.
    \end{example*}
    
    \begin{proof}[Proof of theorem]
        Let $S = \Stab_G(x)$, and let our as yet not well defined map $\Phi$ be such that $\Phi(gS) = gx$.
        \begin{itemize}
            \item \emph{Well-definedness}. For any $g_1$, $g_2 \in G$ with $g_1S = g_2S$, we must have $\Phi(g_1S) = \Phi(g_2S)$, i.e. $g_1x = g_2x$. Since $g_1S = g_2S$, we have $s \in S$ with $g_1 = g_2s$. Then $g_1x = (g_2s)x = g_2(sx) = g_2x$ as required.
            \item \emph{Surjectivity}. For any $gx \in Gx$, we clearly have $\Phi(gS) = gx$.
            \item \emph{Injectivity}. Suppose $\Phi(g_1S) = \Phi(g_2S)$, that is $g_1x = g_2x$. Let $s = g_2^{-1}g_1$, so $sx = (g_2^{-1}g_1)x = x$, and so $s \in S$. But $g_1 = g_2s$, so $g_1 \in g_2S$, and so $g_1S = g_2S$ as required (since cosets partition $G$).
        \end{itemize}
    \end{proof}
    
    \begin{example*}[Symmetries of a cube]
        Let $G$ be the group of isometries of a cube, and let $x$ be the centre of a face. Clearly $|Gx| = 6$, since there are 6 other faces to take $x$ to, and $\Stab_G(x)$ must fix the face on which $x$ lies, so it is isomorphic $D_8$. Thus $|\Stab_G(x)| = 8$, and so $|G| = 48$. Nice!
    \end{example*}
    
    \lecture{Lecture 11 -- 03/11/25}

    The next theorem is another cool application of the orbit-stabiliser theorem.

    \begin{theorem}[Cauchy]
        If $G$ is a finite group and $p$ is a prime dividing $|G|$, there exists $g \in G$ with $o(g) = p$.
    \end{theorem}
    
    \begin{proof}
        Consider the following set $X$ consisting of $p$-tuples of elements of $G$:
        \[
            X = \{(g_1, \dots, g_p): g_i \in G, g_1 \cdots g_p = e\}
        \]
        Define an action of $C_p$ on $X$: if $C_p = \{1, t, t^2, \dots, t^{p-1}\}$, let
        \[
            t^k(g_1, \dots, g_p) = (g_{k+1}, \dots, g_p, g_1, \dots, g_k)
        \]
        in other words, it cyclically permutes the tuple. Let us now verify that this is a bona fide action. It's easy to see that $t^k \cdot t^l (g_1, \dots, g_p) = t^{k+l} (g_1, \dots, g_p)$, and so it remains to verify
        \[
            (g_{k+1}, \dots, g_p, g_1, \dots, g_k) \in X
        \]
        Let us then suppose $g_1 \cdots g_p = e$. For convenience, write
        \begin{align*}
            a &= g_1 \cdots g_k \\
            b &= g_{k+1} \cdots g_p
        \end{align*}
        We know that $a \cdot b = e$, so $b = a^{-1}$ and hence $b \cdot a = e$ -- $a$ and $b$ commute. From this it follows that
        \[
            g_{k+1} \dots g_p \cdot g_1 \dots g_k = e
        \]
        as required. Let us now compute $|X|$. Observe that, for any choices $g_1, \dots, g_{p-1} \in G$, $g_p = g_{p-1}^{-1} \cdots g_1^{-1}$ is the unique choice for $g_n$ such that $(g_1, \dots, g_p) \in X$, and so $|X| = |G|^{p-1}$. Furthermore, the action of $C_p$ partitions $X$ into orbits:
        \[
            X = C_p x_1 \cup C_p x_2 \cup \cdots \cup C_p x_k
        \]
        Then, by \cref{orbit-stabiliser-theorem} (orbit-stabiliser), for each $j$,
        \[
            |C_p x_j| \;\big\vert\; |C_p| = p
        \] 
        so either $|C_p x_j| = 1$ or $|C_p x_j| = p$. Let $l$ be the number of orbits of size 1, so after renumbering we may assume
        \[
            |C_p x_j| =
            \begin{cases}
                1 & 1 \leq j \leq l \\
                p & l < j \leq k
            \end{cases}
        \]
        Since the orbits partition $X$,
        \begin{align*}
            |G|^{p-1} &= |X| \\
            &= \sum_{j=1}^k |C_px_j| \\
            &= l + p(k-l)
        \end{align*}
        and so, since $p \mid |G|$, $p \mid l$. But note that $x = (g_1, \dots, g_p)$ has an orbit of size 1 if and only if
        \[
            g_1 = \cdots = g_p
        \]
        and we can observe that $(e, \dots, e)$ is clearly in $X$ and has an orbit of size 1, so $l > 0$. But then $l \geq p > 1$, so there must be a further trivial orbit
        \[
            (g, \dots, g) \in X
        \]
        with $g \neq e$ and $g^p = e$. Thus $o(g) \mid p$, and so $o(g) = p$ as required.
        
    \end{proof}

    \subsection{Conjugation}

    \begin{definition}
        Let $G$ be a group and $g \in G$, For any $h \in G$,
        \[
            hgh^{-1}
        \]
        is the \emph{conjugate} of $g$ by $h$.
    \end{definition}
    Intuition:

    \begin{itemize}
        \item $hgh^{-1}$ has the `same shape' as $g$;
        \item $hgh^{-1}$ corresponds to `changing the coordinates' of $g$.
    \end{itemize}

    \begin{definition}
        The \emph{conjugacy class} of $g \in G$ is
        \[
            \ccl(g) = \{hgh^{-1} \mid g \in G\}
        \]
    \end{definition}
    
    \begin{example*}
        \begin{itemize}
            \item If $G$ is abelian, then $hgh^{-1} = ghh^{-1} = g$, so the only conjugate of $g$ is $g$ itself.
            \item $D_{2n}$ has one conjugacy class of reflections if $n$ is odd and two if $n$ is even.
        \end{itemize}
    \end{example*}
    
    $G$ acts on itself by conjugation, where we define $g\gamma = g\gamma g^{-1}$. This is an important action which is \emph{very} different to the regular action. Notably, $\ccl(g)$ is the orbit of $g$ under this action.

    \begin{definition}
        The \emph{centraliser} of $g \in G$ is defined as
        \[
            C_G(g) = \{h \in G: hgh^{-1} = g \iff hg = gh\}
        \]
        which is the stabiliser of $g$ under the conjugation action. It is also clearly the set of elements that commute with $g$.
    \end{definition}
    
    \begin{definition}
        The \emph{centre} of $G$ is
        \begin{align*}
            Z(G) &= \{h \in G: hg = gh \text{ for all } g \in G\} \\
            &= \bigcap_{g \in G} C_G(g)
        \end{align*}
        the set of elements that commute with all elements of $G$.
    \end{definition}
    
    
    
    

    \lecture{Lecture 12 -- 05/11/25}

    \section{The M\"obius group}

    This group is \emph{almost} a group of bijections of $\CC$, \emph{but} we need to add a formal `point at $\infty$'.

    \begin{definition}
        The \emph{Riemann sphere} $\CC_\infty$ is $\CC \cup \{\infty\}$.
    \end{definition}
    
    \begin{remark*}
        Why a sphere? We see there is a map $\pi: \CC \to$ a sphere $S^2$, called \emph{stereographic projection}, which misses out the north pole. Hence, by adding $\infty$ to $\CC$ and mapping this to the north pole, we have a bijection between $\CC_\infty$ and $S^2$.
        DIAGRAM!
    \end{remark*}
    
    \begin{definition}
        Let $a$, $b$, $c$, $d \in \CC$, with $ad - bc \neq 0$. If $c \neq 0$, then the corresponding \emph{M\"obius transformation}
        \begin{align*}
            \mu: \CC_\infty &\to \CC_\infty \\
            z &\mapsto \begin{cases}
                \displaystyle \frac{az+b}{cz+d} & \displaystyle z \neq -\frac{d}{c} \\[10pt]
                \infty & \displaystyle z = -\frac{d}{c} \\[10pt]
                \displaystyle \frac{a}{c} & z = \infty
            \end{cases}
        \end{align*}
        If $c = 0$, then
        \begin{align*}
            \mu: \CC_\infty &\to \CC_\infty \\
            z &\mapsto \begin{cases}
                \displaystyle \frac{az+b}{d} & \displaystyle z \neq -\frac{d}{c} \\[10pt]
                \infty & z = \infty
            \end{cases}
        \end{align*}
        But we usually just write
        \[
            \mu(z) = \frac{az+b}{cz+d}
        \]
        and all the edge cases are implicit.

        The \emph{M\"obius} group $\mathcal M$ is then the group of M\"obius transformations under composition.
    \end{definition}
    
    \begin{proposition}
        If
        \[
            \mu_1(z) = \frac{a_1 z + b_1}{c_1 z + d_1}, \quad \mu_2(z) = \frac{a_2 z + b_2}{c_2 z + d_2}
        \]
        then
        \[
            \mu_1 \circ \mu_2 (z) = \frac{(a_1a_2 + b_1c_2)z + (a_1b_2 + b_1d_2)}{(c_1a_2 + d_1c_2)z + (c_1b_2 + d_1d_2)}
        \]
        \label{mobius-composition}
    \end{proposition}

    \begin{proof}
        Exercise.
    \end{proof}    
    
    \begin{remark*}
        Compare
        \[
            \begin{pmatrix}
             a_1 & b_1 \\
             c_1 & d_1 \\
            \end{pmatrix}\begin{pmatrix}
             a_2 & b_2 \\
             c_2 & d_2 \\
            \end{pmatrix} = \begin{pmatrix}
             a_1a_2 + b_1c_2 & a_1b_2 + b_1d_2 \\
             c_1a_2 + d_1c_2 & c_1b_2 + d_1d_2 \\
            \end{pmatrix}
        \]
        so the coefficients in the composition are given by matrix multiplication.
    \end{remark*}

    \begin{theorem}
        $(\mathcal M, \circ, \id)$ is a group.
    \end{theorem}
    
    \begin{proof}
        Composition of functions is associative, so that's ok. The identity is given by
        \[
            \id: z \mapsto z = \frac{1 \cdot z + 0}{0 \cdot z + 1}
        \]
        so that's ok too. Now let's think about inverses. For a M\"obius transformation
        \[
            \mu: z \mapsto \frac{az+b}{cz+d}
        \]
        let
        \[
            \nu: z \mapsto \frac{dz-b}{-cz+a}
        \]
        (whose coefficients are given by the matrix inverse of those of $\mu$). It follows from \cref{mobius-composition} that $\mu \circ \nu = \id$ (here we are using $ad - bc \neq 0$). Closure again follows from \cref{mobius-composition} and the comparison with matrix multiplication.
    \end{proof}
    
    One important way to study M\"obius transformations is via their fixed points.

    \begin{definition}
        Suppose $f: X \to X$ is a permutation. If $f(x) = x$ for some $x \in X$, then $x$ is called a \emph{fixed point} of $f$.
    \end{definition}
    
    \begin{lemma}[3-point lemma for M\"obius transformations]
        If $\mu \in \mathcal M$ fixes three distinct points $w_1$, $w_2$, $w_3 \in \CC_\infty$, then $\mu = \id$.
    \end{lemma}
    
    \begin{proof}
        Let
        \[
            \mu(z) = \frac{az+b}{cz+d}
        \]
        Then, a fixed point $w_i$ satisfies
        \[
            w_i = \frac{aw_i + b}{cw_i + d}
        \]
        If $w_1 = \infty$, then $c = 0$, so $w_2$ and $w_3$ satisfy
        \begin{align*}
            w_i &= \frac{aw_i + b}{d} \\
            \iff (a-d)w_i + b &= 0
        \end{align*}
        which has two distinct roots $w_2$ and $w_3$. It follows that $a = d$, $b = 0$. Then it's clear that $\mu = \id$.

        Now suppose that none of the $w_i$ are infinity. Again,
        \begin{align*}
            w_i &= \frac{aw_i + b}{cw_i + d} \\
            \implies cw_i^2 + dw_i &= aw_i + b \\
            \implies cw_i^2 + (d-a)w_i - b &= 0
        \end{align*}
        Now we have a quadratic with 3 distinct roots, so again, we must have $b = c = 0$, $a = d$. As before, $\mu = \id$.
    \end{proof}
    
    \begin{proposition}
        Every $\mu \in \mathcal M$ has at least one fixed point.
    \end{proposition}

    \begin{proof}[Sketch proof]
        As above -- every quadratic has at least one root in $\CC$. Full proof an exercise
    \end{proof}
    
    \begin{example*}
        \begin{itemize}
            \item $\mu(z) = z+1$ has $\infty$ as its unique fixed point -- $\Fix(\mu) = \{\infty\}$.
            \item $\mu(z) = 2z$ has 0 and $\infty$ as its two fixed points -- $\Fix(\mu) = \{0, \infty\}$.
        \end{itemize}
    \end{example*}
    
    \begin{lemma}[Triple transitivity]
        For any two triples of distinct points $z_1$, $z_2$, $z_3$ and $w_1$, $w_2$, $w_3 \in \CC_\infty$, there exists $\mu \in \mathcal M$ such that $\mu(z_i) = w_i$ for each $i$.
    \end{lemma}
    
    \begin{proof}
        Let
        \[
            \alpha(z) = \frac{z_2-z_3}{z_2-z_1} \cdot \frac{z-z_1}{z-z_3}
        \]
        (modify appropriately if any $z_i = \infty$). Note
        \begin{align*}
            z_1 &\longmapsto 0 \\
            z_2 &\longmapsto \infty \\
            z_3 &\longmapsto 1
        \end{align*}
        Similarly, we can write down $\beta$ which sends $w_1$, $w_2$, $w_3$ to these three points. Then, $\beta^{-1} \circ \alpha$ is the M\"obius transformation we're after.
        
    \end{proof}
    
    \begin{remark*}
        Named for the fact that it's a strengthening of transitivity when interpreting $\mathcal M$ as acting on $\CC_\infty$.
    \end{remark*}
    
    \begin{remark*}
        The three-point lemma implies that this $\mu$ is unique. We then say further that the action $\mathcal M \acts \CC_\infty$ is \emph{sharply triply transitive}.
    \end{remark*}
    
    \begin{definition}
        Let $z_1$, $z_2$, $z_3$, $z_4 \in \CC_\infty$ be distinct. Since $\mathcal M \acts \CC_\infty$ sharply triply transitively, there is a unique $\alpha$ such that $\alpha(z_1) = 0$, $\alpha(z_2) = 1$, and $\alpha(z_3) = \infty$. Then the \emph{cross ratio}
        \[
            [z_1, z_2, z_3, z_4] = \alpha(z_4)
        \]
        If we really want, we can find an explicit formula using our explicit formula for $\alpha$:
        \[
            [z_1, z_2, z_3, z_4] = \frac{z_2-z_3}{z_2-z_1} \cdot \frac{z_4-z_1}{z_4-z_3}
        \]
    \end{definition}
    
    \lecture{Lecture 13 -- 07/11/25}

    Let's try and find a generating set for the M\"obius group.
    
    \begin{proposition}[Generating set for the M\"obius group]
        $\mathcal M$ is generated by the set of elements of the following 3 forms:
        \begin{itemize}
            \item $\alpha_a: z \mapsto az$, $a \neq 0$;
            \item $\beta_b: z \mapsto z + b$, $b \in \CC$;
            \item $\gamma: z \to \frac{1}{z}$.
        \end{itemize}
        \label{generators-mobius-group}
    \end{proposition}
    
    \begin{proof}
        Let $\mu \in \mathcal M$ be arbitrary, and let $z_1 = \mu(0)$, $z_2 = \mu(1)$, $z_3 = \mu(\infty)$.
        \begin{enumerate}[label=\protect\circled{\arabic*}]
            \item We want to construct $\mu_1$ such that $\mu_1(z_3) = \infty$. Either $z_3 = \infty$ and $\mu_1 = \id$, or
            \[
                \mu_1(z) = \frac{1}{z+b} = \gamma \circ \beta_b(z)
            \]
            where $b = -z_3$. Now, let $z_1' = \mu_1(z_1)$, $z_2' = \mu_1(z_2)$.
            \item Let $b' = -z_1'$ and $\mu_2 = \beta_b'$, that is $\mu_2(z) = z + b'$. Observe $\mu_2(\infty) = \infty$ and $\mu_2(z_1') = 0$. Then by constructikn we have
            \begin{align*}
                \mu_2 \circ \mu_1 (z_3) &= \infty \\
                \mu_2 \circ \mu_1 (z_1) &= \mu_2(z_1') = 0
            \end{align*}
            Lastly, let $z_2'' = \mu_2 \circ \mu_1(z_2) \neq 0$, $\infty$.
            \item Let $a = \frac{1}{z_2''}$ and $\mu_3(z) = \alpha_a(z)$, that is $\frac{z}{z_2''}$.
        \end{enumerate}
        It's now clear that 0, 1 and $\infty$ have the same image under $\mu$ and $\mu_1^{-1} \circ \mu_2^{-1} \circ \mu_3^{-1}$, from which it follows by the three-point lemma that $\mu = \mu_1^{-1} \circ \mu_2^{-1} \circ \mu_3^{-1}$. Hence our generating set generates $\mu$ as required.
    \end{proof}
    
    \begin{definition}
        A \emph{circle} in $\CC_\infty$ is
        \begin{itemize}
            \item a standard Euclidean circle in $\CC$;
            \item or $l \cup \{\infty\}$ where $l \subseteq \CC$ is a Euclidean straight line.
        \end{itemize}
    \end{definition}
    
    \begin{remark}
        For completeness, let us add that Euclidean circles in $\CC$ have the form $|z-c| = r$, $c \in \CC$, $r > 0 \in \RR$, while Euclidean lines have the form $|z-a| = |z-b|$ for $a$, $b$ not unique in $\CC$.
    \end{remark}
    
    \begin{theorem}[Circles and M\"obius transformations]
        If $C \subseteq \CC_\infty$ is a circle and $\mu \in \mathcal M$ is a M\"obius transformation, then $\mu(C)$ is also a circle.
    \end{theorem}
    
    \begin{proof}
        This is a great opportunity to apply \cref{generators-mobius-group} -- it is sufficient to check the result for only $\alpha_a(z) = az$, $\beta_b(z) = z + b$, $\gamma(z) = \frac{1}{z}$. Furthermore, it is really quite obvious that $\alpha_a$ and $\beta_b$ preserve circles.

        If $C$ is a Euclidean circle, $\gamma(C)$ has equation
        \begin{align*}
            \left|\frac{1}{z} - c\right| &= r \\
            \iff \left(\frac{1}{z} - c\right)\left(\frac{1}{\bar z} - \bar c\right) &= r^2 \\
            \iff \frac{1}{|z|^2} - \frac{c}{z} - \frac{\bar c}{z} + |c|^2 - r^2 &= 0 \\
            \iff (|c|^2 - r^2)|z|^2 - cz - \bar c \bar z + 1 &= 0
        \end{align*}
        In the case $|c|^2 = r^2$, we have
        \begin{align*}
            cz + \bar c \bar z &= 1 \\
            \iff \frac{z}{\bar c} + \frac{\bar z}{c} &= \frac{1}{|c|^2} \\
            \iff |z|^2 &= |z|^2 - \frac{z}{\bar c} - \frac{\bar z}{c} + \frac{1}{|c|^2} \\
            &= \left|z-\frac{1}{c}\right|^2 \\
            \iff |z| &= \left|z - \frac{1}{c}\right|
        \end{align*}
        which is a straight line -- furthermore, it is clear 0 is on the circle, and this point will go to infinity. Otherwise, $|c|^2 \neq r^2$, so we have
        \begin{align*}
            |z|^2 - \left(\frac{c}{|c|^2 - r^2}\right)z - \left(\frac{\bar c}{|c|^2 - r^2}\right) \bar z + \frac{1}{|c|^2 - r^2} &= 0 \\
            \iff \left|z - \frac{\bar c}{|c|^2 - r^2}\right|^2 &= \frac{|c|^2}{(|c|^2 - r^2)^2} - \frac{1}{|c|^2 - r^2} \\
            &= \frac{r^2}{(|c|^2 - r^2)^2}
        \end{align*}
        which is a circle. The case where $C$ is a line union infinity is left as an exercise.        
    \end{proof}
    
    \begin{corollary}
        Four points $z_1$, $z_2$, $z_3$, $z_4 \in \CC$ lie on a circle if and only if
        \[
            [z_1, z_2, z_3, z_4] \in \RR \cup \{\infty\}
        \]
    \end{corollary}
    
    \begin{proof}
        Let $\alpha \in \mathcal M$ be such that $\alpha(z_1) = 0$, $\alpha(z_2) = 1$, $\alpha(z_3) = \infty$, and $\alpha(z_4) = [z_1, z_2, z_3, z_4]$.
        \begin{itemize}
            \item[$\implies$] Suppose $z_1, \dots, z_4 \in C$ a circle. Then 0, 1, $\infty$ and $\alpha(z_4)$ will lie on $\alpha(C)$, also a circle. Since $\infty$ is on this circle, it is in fact a line union infinity, and so, since it passes through $0$ and $1$, it is $\RR \cup \{\infty\}$, so $\alpha(z_4) \in \RR \cup \{\infty\}$ as required.
            \item[$\impliedby$] Suppose $\alpha(z_4) \in \RR \cup \{\infty\}$, and of course $\alpha(z_1) = 0$, $\alpha(z_2) = 1$, $\alpha(z_3) = \infty$ -- so they all lie on a circle. Hence $z_1, \dots, z_4 \in \alpha^{-1}(\RR \cup \{\infty\})$ which is a circle by the theorem. 
        \end{itemize} 
    \end{proof}
    
    \section{Finite groups}
    
    Let's categorise all the finite groups!

    \lecture{Lecture 14 -- 10/11/25}

    \begin{itemize}
        \item \emph{Order 1}. If $|G| = 1$, then $G \cong \{e\}$, the trivial group.
        \item \emph{Order 2}. If $|G| = 2$, since 2 is prime, $G \cong C_2$ by \cref{prime-groups-cyclic}.
        \item \emph{Order 3}. As above, $G \cong C_3$.
        \item \emph{Order 4}. If $|G| = 4$, we have $C_4$. But is that all?
    \end{itemize}

    \begin{definition}[Direct product]
        If $G$ and $H$ are groups, the \emph{direct product} is
        \[
            G \times H = \{(g, h): g \in G, h \in H\}
        \]
        with group operation
        \[
            (g_1, h_1) \cdot (g_2, h_2) = (g_1g_2, h_1h_2)
        \]
        and identity $(e_G, e_H)$. Inverses are also pointwise -- $(g, h)^{-1} = (g^{-1}, h^{-1})$.
    \end{definition}
    
    \begin{example*}[Klein 4-group]
        The group
        \[
            K_4 := C_2 \times C_2
        \]
        is often called the Klein 4-group (it also makes sense to think of $D_4$ as $K_4$). Note that every element in $K_4$ has order 2, while there exist elements of order 4 in $C_4$, so $K_4 \not\cong C_4$.
    \end{example*}

    \begin{theorem}[Direct product theorem]
        If $H_1$, $H_2 \leq G$ and
        \begin{itemize}
            \item $H_1 \cap H_2 = \{e\}$;
            \item $h_1h_2 = h_2h_1$ for all $h_1 \in H_1$, $h_2 \in H_2$;
            \item $G = H_1H_2$ -- for all $g \in G$, $g = h_1h_2$ for some $h_1 \in H_1$, $h_2 \in H_2$.
        \end{itemize}
        Then $G \cong H_1 \times H_2$.
        \label{direct-product-theorem}
    \end{theorem}
    
    \begin{proof}
        Define
        \begin{align*}
            \Phi: H_1 \times H_2 &\to G \\
            (h_1, h_2) &\mapsto h_1h_2
        \end{align*}
        We claim $\Phi$ is an isomorphism.
        \begin{itemize}
            \item \emph{$\Phi$ is a homomorphism}. For any $h_1$, $h_1' \in H_1$, $h_2$, $h_2' \in H_2$, we have
            \begin{align*}
                \Phi(h_1, h_2)\Phi(h_1', h_2') &= h_1h_2h_1'h_2' \\
                &= h_1h_1'h_2h_2' \\
                &= \Phi(h_1h_1', h_2h_2') \\
                &= \Phi((h_1, h_2)(h_1', h_2'))
            \end{align*}
            \item \emph{$\Phi$ is surjective}. Follows immediately from $G = H_1H_2$.
            \item \emph{$\Phi$ is injective}. We will show $\ker \Phi$ is trivial. Suppose $\Phi(h_1, h_2) = e$ for some $(h_1, h_2) \in H_1 \times H_2$. Then
            \begin{align*}
                h_1h_2 &= e \\
                \implies H_1 \ni h_1 &= h_2^{-1} \in H_2
            \end{align*}
            so $h_1 = h_2^{-1} \in H_1 \cap H_2$. But $H_1 \cap H_2 = \{e\}$, so $h_1 = h_2 = e$. Thus $(h_1, h_2)$ is the identity in $H_1 \times H_2$ as required.            
        \end{itemize}
    \end{proof}
    
    \begin{remark*}
        If we know $H_1 \cap H_2 = \{e\}$, then $|H_1 \times H_2| = |H_1||H_2|$. So if $|H_1||H_2| = |G|$, then we get that $G = H_1H_2$ automatically.
    \end{remark*}
    
    \begin{lemma}[Groups of order 4]
        If $|G| = 4$, then either $G \cong C_4$ or $G \cong K_4$.
    \end{lemma}
    
    \begin{proof}
        By Lagrange, every non-trivial element has order 2 or 4. If there is an element $g \in G$ of order 4, then $G = \gen g \cong C_4$. Otherwise, every non-trivial element of $G$ has order 2. Let $a$, $b$ be distinct non-trivial elements of $G$, and let $H_1 = \gen a$, $H_2 = \gen b$. Clearly $H_1 \cap H_2 = \{e\}$ and, since $|H_1| = |H_2| = 2$, we have $G = H_1H_2$. Finally, since $o(ab) = 2$, $abab = e \implies ab = b^{-1}a^{-1} = ba$, so we have commutativity. Thus $G = H_1 \times H_2 \cong C_2 \times C_2$.
    \end{proof}
    
    \begin{theorem}[Chinese remainder theorem]
        If $\gcd(m, n) = 1$, then $C_m \times C_n \cong C_{mn}$.
        \label{chinese-remainder-theorem}
    \end{theorem}
    
    \begin{proof}
        Let $C_{mn} = \gen g$. Let $H_1 = \gen{g^n} \cong C_m$, $H_2 = \gen{g^m} \cong C_n$. We will check the hypotheses of the direct product theorem.
        \begin{itemize}
            \item Observe $g^k \in H_1 \iff k = np + mnq$ for some $p$, $q \in \ZZ$, i.e. if and only if $n \mid k$. Similarly, $g^k \in H_2$ if and only if $m \mid k$. Hence $g^k \in H_1 \cap H_2$ if and only if $n, m \mid k \iff \lcm(m, n) \mid k$. But $\gcd(m, n) = 1 \implies \lcm(m, n) = mn$, so $mn \mid k$ -- but $g^{mn} = e$, so $g^k = e$ as required.
            \item $C_{mn}$ is abelian, so we definitely have commutativity.
            \item Immediate by the remark.
        \end{itemize}
    \end{proof}
    
    Let's continue classifying groups!

    \begin{itemize}
        \item \emph{Order 5}. $C_5$ is of course the only possibility.
        \item \emph{Order 6}.
    \end{itemize}
    
    \begin{lemma}[Groups of order 6]
        If $|G| = 6$, then $G \cong C_6$ or $G \cong D_6$.
    \end{lemma}
    
    \begin{proof}
        By Cauchy's theorem, $G$ has an element $r$ of order 3 and an element $s$ of order 2. Since $|\gen r| = 3$, $|G:\gen r| = 2$, and $s \notin \gen r$, so $s \gen r = G \setminus \gen r = \gen r s$. We conclude that $sr = r^is$ for some $i = 0, 1, 2$.

        \emph{Case 1}. $i = 0$. Then $sr = s \implies r = e$, but this contradicts $o(r) = 3$.

        \emph{Case 2}. $i = 1$. Then $sr = rs$, from which it's clear that $\gen r$ and $\gen s$ commute. Then \cref{direct-product-theorem} (the direct product theorem) applies, so $G \cong C_2 \times C_3 = C_6$ by \cref{chinese-remainder-theorem}.

        \emph{Case 3}. Then we have $sr = r^2s = r^{-1}s$ -- this is the dihedral relation. We conclude $G \cong D_6$ by proposition 1.31 (to be added).
    \end{proof}
    
    \begin{remark*}
        $S_3$ is a non-abelian group of order 6, so this gives us a way to conclude $S_3 \cong D_6$.
    \end{remark*}

    \lecture{Lecture 15 -- 12/11/25}
    
    \begin{itemize}
        \item \emph{Order 7}. $C_7$ only.
        \item \emph{Order 8}. Interesting! We have $C_2 \times C_2 \times C_2$, $C_4 \times C_2$, $C_8$ of course. We also have $D_8$. But there's another one!
    \end{itemize}
    
    \begin{example}[The quaternion group]
        Let
        \begin{align*}
            Q_8 = \left\{\begin{pmatrix}
             \pm 1 & 0 \\
             0 & \pm 1 \\
            \end{pmatrix}, \begin{pmatrix}
             \pm i & 0 \\
             0 & \mp i \\
            \end{pmatrix}, \begin{pmatrix}
             0 & \pm 1 \\
             \mp 1 & 0 \\
            \end{pmatrix}, \begin{pmatrix}
             0 & \pm i \\
             \pm i & 0 \\
            \end{pmatrix}\right\}
        \end{align*}
        One can check that this is a group. We usually use Hamilton's notation, where
        \begin{align*}
            1 = \begin{pmatrix}
             1 & 0 \\
             0 & 1 \\
            \end{pmatrix},\;
            i = \begin{pmatrix}
             i & 0 \\
             0 & -i \\
            \end{pmatrix},\;
            j = \begin{pmatrix}
             0 & 1 \\
             -1 & 0 \\
            \end{pmatrix},\;
            k = \begin{pmatrix}
             0 & i \\
             i & 0 \\
            \end{pmatrix}
        \end{align*}
        so that the group consists of $\{\pm 1, \pm i, \pm j, \pm k\}$. Then we have various properties:
        \begin{itemize}
            \item $i^2 = j^2 = k^2 = -1$;
            \item $(-1)i = -i$, etc;
            \item $ij = k$, $jk = i$, $ki = j$;
            \item $ji = -k$, $kj = -i$, $ik = -j$.
            \item $-1$ commutes with everything.
        \end{itemize}
        But is this in our list? It's not abelian, so if it were, it would have to be $D_8$. But $D_8$ has too many elements of order 2.
    \end{example}
    
    \begin{lemma}[Groups of order 8]
        If $|G| = 8$, then $G \cong C_2 \times C_2 \times C_2$, $C_4 \times C_2$, $C_8$, $D_8$ or $Q_8$.
    \end{lemma}
    
    \begin{proof}
        By Lagrange, every element has order 1, 2, 4 or 8.
        \begin{itemize}
            \item If we have an element of order 8, we're done ($G \cong C_8$).
            \item Suppose that every non-trivial element has order 2. We know it's abelian, and then by picking $a$, $b \neq a$, $c \neq a, b, ab$, we see that
            \[
                G = \gen a \times \gen b \times \gen c \cong C_2 \times C_2 \times C_2
            \]
            by applying the direct product theorem twice.
            \item Now we may suppose we have $a \in G$ with $o(a) = 4$, so $|G: \gen a| = 2$ -- also let $b \in G \setminus \gen a$. Then
            \[
                b \gen a = G \setminus \gen a = \gen a b
            \]
            so that $ba = a^ib$ for some $i = 0, 1, 2, 3$.
            \begin{itemize}
                \item $i = 0$, so $a = e$, which is a contradiction.
                \item $i = 1$, so $ba = ab$, and it follows $ba^j = a^jb$ for all $j$. It follows from here that the whole group is abelian.
                \item $i = 2$, so $ba = a^2b \implies bab^{-1} = a^2$. But $bab^{-1}$ is $a$ conjugated by $b$, so must have the same order as $a$, i.e. 4, while $a^2$ clearly has order 2. So this is a contradiction.
                \item $i = 3$, so $ba = a^3 b = a^{-1} b$.
            \end{itemize}
            Let's take a break and think about powers of $b$. If $b^2 = ba^j$, then $b = a^j$, so $b \in \gen a$, but we picked $b \in G \setminus \gen a$. So $b^2 \in \gen a$, and we get four more cases.
            \begin{itemize}
                \item $b^2 = e$.
                \item $b^2 = a$. Then, since $a$ has order 4, $b$ has order 8, so we get $C_8$ again.
                \item $b^2 = a^2$.
                \item $b^2 = a^3 = a^{-1}$ -- again, $a^{-1}$ has order 4, and so we get $C_8$ once more.
            \end{itemize}
            This leaves us with four combinations of cases to analyse.
            \begin{itemize}
                \item $G$ is abelian, $b^2 = e$. Then $\gen b \cong C_2$, $\gen a \cong C_4$, $\gen b \cap \gen a = \{e\}$, so $G \cong C_4 \times C_2$ by the direct product theorem.
                \item $G$ is abelian, $b^2 = a^2$. Then $ab^{-1}$ has order 2, and it turns out that once again $G \cong C_4 \times C_2$ by the direct product theorem.
                \item $ba = a^{-1}b$, $b^2 = e$, and we already knew $a^4 = e$. That's enough to show that $G \cong D_8$ with $r = a$ and $s = b$.
                \item $ba = a^{-1}b$, $b^2 = a^2$. We claim $G \cong Q_8$. Let $i = a$, $j = b$, $k = ab$, $-1 = a^2 = b^2$. Then
                \begin{align*}
                    G &= \{e, a, a^2, a^3, b, ab, a^2b, a^3b\} \\
                    &= \{1, i, -1, -i, j, k, -j, -k\}
                \end{align*}
                and we can check multiplication works. So $G \cong Q_8$ as required.
            \end{itemize}
        \end{itemize}
    \end{proof}
    
    We have now categorised all the groups of order less than or equal to 8: $\{e\}$, $C_2$, $C_3$, $C_4$, $C_2 \times C_2$, $C_5$, $C_6$, $D_6$, $C_7$, $C_8$, $C_4 \times C_2$, $C_2 \times C_2 \times C_2$, $D_8$, $Q_8$.

    \section{Quotients}
    
    \subsection{Normal subgroups}

    \begin{definition}
        $H \leq G$ is called \emph{normal} if $ghg^{-1} \in H$ for every $h \in H$ and $g \in G$. We write $H \triangleleft G$.
    \end{definition}
    
    \begin{example}
        \begin{itemize}
            \item $\{e\} \triangleleft G$, $G \triangleleft G$ for all $G$;
            \item If $G$ is abelian, $H \triangleleft G$ for all $H \leq G$;
            \item $\gen r \triangleleft D_{2n}$, but $\gen s$ is not a normal subgroup of $D_{2n}$.
        \end{itemize}
    \end{example}
    
    \begin{proposition}
        Let $\phi: G \to H$ be a homomorphism. Then $\ker \phi \triangleleft G$.
    \end{proposition}
    
    \begin{proof}
        We have
        \begin{align*}
            \phi(ghg^{-1}) &= \phi(g)\phi(h)\phi(g)^{-1} \\
            &= \phi(g)\phi(g)^{-1} \\
            &= e
        \end{align*}
        so $ghg^{-1} \in \ker \phi$. Hence $\ker \phi \triangleleft G$ as required.
    \end{proof}

    \lecture{Lecture 16 -- 14/11/25}
    
    \begin{lemma}
        Suppose $H \leq G$. Then $H \triangleleft G$ if and only if $gH = Hg$ for all $g \in G$.
        \label{normal-subgroups-equal-cosets}
    \end{lemma}
    
    \begin{proof}
        \begin{itemize}
            \item[$\implies$] Let $h \in H$, $g \in G$. Since $H \triangleleft G$, $ghg^{-1} \in H$. Hence $gh = (ghg^{-1})g \in Hg$, so $gH \subseteq Hg$. Similarly, $Hg \subseteq gH$, so $gH = Hg$ as required.
            \item[$\impliedby$] Suppose $gH = Hg$. Let $h \in H$, so $gh \in gH = Hg$, so we have $h' \in H$ with $gh = h'g$. Hence $ghg^{-1} = h' \in H$ as required. 
        \end{itemize}
    \end{proof}
    
    \begin{theorem}[Quotient groups]
        If $H \triangleleft G$, the set of left cosets $G/H$ is a group with operation $(g_1H)(g_2H) = (g_1g_2)H$.
        \label{quotient-is-a-group}
    \end{theorem}
    
    \begin{proof}
        We need to check that the operation is well-defined and satisfies the group axioms.
        \begin{itemize}
            \item \emph{Well-defined}. Suppose $g_1H = g_1'H$ and $g_2H = g_2'H$, so $Hg_2 = Hg_2'$ (since $H$ normal in $G$). Hence, we have $h_1$, $h_2 \in H$ with $g_1 = g_1'h_1$ and $g_2 = h_2g_2'$. Therefore,
            \begin{align*}
                g_1g_2 &= (g_1'h_1)(h_2g_2') \\
                &= g_1'(h_1h_2)g_2' \\
                &= g_1'g_2'h_3
            \end{align*}
            for some $h_3 \in H$ by \cref{normal-subgroups-equal-cosets}. Therefore, $g_1g_2 \in g_1'g_2'H$, so $g_1g_2H = g_1'g_2'H$.
            \item \emph{Associativity}. Obvious from the definition.
            \item \emph{Identity}. $H$.
            \item \emph{Inverses}. Clearly $g^{-1}H = (gH)^{-1}$.
        \end{itemize}
    \end{proof}
    
    \begin{definition}
        If $H \triangleleft G$, then $G/H$, which is a group by \cref{quotient-is-a-group}, is called the \emph{quotient of $G$ by $H$}.
    \end{definition}
    
    \begin{example}
        \begin{itemize}
            \item $G/\{e\} \cong G$, $G/G \cong \{e\}$.
            \item Since $\ZZ$ is abelian, $n\ZZ \triangleleft \ZZ$ for any $n$, and $\ZZ/n\ZZ$ is cyclic of order $n$, with generator $1 + n\ZZ$, so $\ZZ/n\ZZ \cong C_n$.
            \item Let $G$ be a group, $H \leq G$, and suppose $|G:H| = 2$. Then $H \triangleleft G$ as we've argued before. For example, $C_n \cong \gen r \triangleleft D_{2n}$, and $D_{2n} / C_n \cong C_2$.
        \end{itemize}
    \end{example}
    
    \begin{remark*}
        Quotienting and the direct product are not opposites! $A / B \cong C \nRightarrow A \cong B \times C$.
    \end{remark*}
    
    \begin{theorem}[Isomorphism theorem]
        If $\phi: G \to H$ is a homomorphism, then
        \[
            G/\ker \phi \cong \im \phi
        \]
        \label{isomorphism-theorem}
    \end{theorem}
    
    \begin{proof}
        Since $\ker \phi \triangleleft G$, the quotient $G/\ker \phi$ is indeed a group. Let's define
        \begin{align*}
            \bar\phi: G/\ker\phi &\to \im \phi \\
            g \ker \phi &\mapsto \phi(g)
        \end{align*}
        We must perform all the usual checks.
        \begin{itemize}
            \item \emph{Well-defined}. Suppose $g_1 \ker \phi = g_2 \ker \phi$, so $g_1 = g_2 k$ for some $k \in \ker \phi$. Then
            \begin{align*}
                \bar\phi(g_1 \ker \phi) &= \phi(g_1) \\
                &= \phi(g_2k) \\
                &= \phi(g_2)\phi(k) \\
                &= \phi(g_2) \\
                &= \bar\phi(g_2 \ker \phi)
            \end{align*}
            \item \emph{Homomorphism}. For $g_1$, $g_2 \in G$,
            \begin{align*}
                \bar\phi((g_1 \ker \phi)(g_2 \ker \phi)) &= \bar\phi(g_1g_2 \ker \phi) \\
                &= \phi(g_1g_2) \\
                &= \phi(g_1)\phi(g_2) \\
                &= \bar\phi(g_1 \ker \phi)\bar\phi(g_2 \ker \phi)
            \end{align*}
            \item \emph{Injective}. If $\bar\phi(g \ker \phi) = e$, then $\phi(g) = \bar\phi(g \ker \phi) = e$, so $g \in \ker \phi$. Hence $g \ker \phi = \ker \phi$, and so $\ker \bar\phi = \{\ker \phi\}$ is trivial, hence $\bar\phi$ is injective as required.
            \item \emph{Surjective}. An element of $\im \phi$ looks like $\phi(g)$, which is $\bar\phi(g \ker \phi)$ as required.           
        \end{itemize}
    \end{proof}
    
    \begin{example}
        \begin{itemize}
            \item Since $\phi: \ZZ \to \CC_\times$, $k \mapsto e^{2\pi i k/n}$ is a homomorphism with image $C_n$ and kernel $n \ZZ$, we get that $\ZZ/n\ZZ \cong C_n$ by the isomorphism theorem.
            \item Similarly, $\phi: \RR \to \CC_\times$, $k \mapsto e^{2\pi i t}$ is a homomorphism with image $\{z: |z| = 1\} = U(1)$ and kernel $\ZZ$, we have $\RR/\ZZ \cong U(1)$.
        \end{itemize}
    \end{example}

    \lecture{Lecture 17 -- 17/11/25}
    
    \begin{definition}
        A group $G$ is \emph{simple} if the only normal subgroups are $\{e\}$ and $G$ itself. Thus every homomorphism out of $G$ is either injective or trivial by the isomorphism theorem.
    \end{definition}
    
    \begin{example}
        $C_p$ is simple for $p$ a prime.
    \end{example}
    
    \begin{remark*}
        The classification of finite simple groups is one of the great mathematical achievements of the $20^\text{th}$ century.
    \end{remark*}
    
    \section{Permutations}

    Recall that a permutation of a set $X$ is a bijection $X \to X$, and that $\Sym(X)$ is the group of all such permutations under composition.

    \begin{definition}
        Any list of $k$ distinct elements $a_1, \dots, a_k \in \{1, \dots, n\}$ determines a \emph{$k$-cycle}:
        \[
            \sigma = (a_1 \; \dots \; a_k)
        \]
        which is a permutation defined as follows:
        \[
            b \mapsto \begin{cases}
                a_{i+1} & b = a_i, i \neq k \\
                a_1 & b = a_k \\
                b & \text{otherwise}
            \end{cases}
        \]
    \end{definition}
    \begin{example}
        The permutation $\sigma$ taking 1 to 3, 2 to 1 and 3 to 2 is a 3-cycle represented as
        \begin{align*}
            \sigma &= (1\; 3\; 2) \\
            &= (3\; 2\; 1) \\
            &= (2\; 1\; 3)
        \end{align*}
        The permutation swapping 1 and 2 and fixing 3 is
        \begin{align*}
            \tau &= (1 \; 2) = (2\; 1)
        \end{align*}
    \end{example}
    Multiplying cycles is not entirely trivial, but one can easily become proficient in the art. The most important thing is to start from the right.
    \begin{example}
        \begin{itemize}
            \item \begin{align*}
                \tau \sigma &= (1 \; 2)(1 \; 3 \; 2) \\
                &= (1\; 3)
            \end{align*}
            \item \[
                (1\; 4\; 3\; 2)(2\; 4\; 3) = (1\; 4\; 2\; 3)
            \]
        \end{itemize}
    \end{example}
    
    \begin{remark*}
        Cycles which are cyclic permutations of one another are the same: $(a_1\; \dots \; a_k) = (a_2 \; \dots \; a_k\; a_1) = \cdots$
    \end{remark*}
    
    \begin{definition}
        Cycles $(a_1 \; \dots \; a_k)$ and $(b_1 \; \dots \; b_l)$ are \emph{disjoint} if they are disjoint as sets -- $a_i \neq b_j$ for any $i$, $j$. Disjoint cycles have the nice property that they commute.
    \end{definition}
    
    \begin{theorem}[Disjoint cycles]
        Every $\sigma \in S_n$ can be written as a product of disjoint cycles. Furthermore, this product is unique up to cyclic permutation of cycle elements and cycle reordering.
    \end{theorem}
    
    \begin{proof}
        The action of $\gen \sigma$ on $X = \{1, \dots, n\}$ partitions $X$ into orbits:
        \[
            X = \gen \sigma i_1 \cup \gen \sigma i_2 \cup \dots \cup \gen \sigma i_k
        \]
        Let $n_j = |\gen \sigma i_j|$. Then
        \[
            \sigma = (i_1 \; \sigma(i_j) \; \dots \; \sigma^{n_1-1}(i_1)) \dots (i_k \; \sigma(i_k) \; \dots \; \sigma^{n_k-1}(i_k))
        \]
        which proves existence. For uniqueness, note that the only choices we made were the orbit representatives in each orbit and the indexing of the orbits/their representatives. These correspond to cyclically permuting the cycles of $\sigma$ and changing the order of the cycles respectively, which are exactly the failures of uniqueness that we are allowed to have.
    \end{proof}
    
    \begin{definition}[Cycle types]
        If $\sigma = (a_1^1 \; \dots \; a_{k_1}^1) \dots (a_1^l \; \dots \; a_{k_l}^l)$ is a product of disjoint cycles then $\sigma$ is called a \emph{$(k_1, \dots, k_l)$-cycle}, and $\{k_1, \dots, k_l\}$ is called $\sigma$'s \emph{cycle type}.
    \end{definition}
    
    \begin{remark*}
        If $\sigma$ is a $k$-cycle, then $\sigma$ has order $k$. Furthermore, if $\sigma$ has cycle type $\{k_1, \dots, k_l\}$, then $\sigma$ has order $\lcm(k_1, \dots, k_l)$.
    \end{remark*}
    
    \begin{definition}
        A \emph{transposition} is just another name for a 2-cycle.
    \end{definition}
    
    \begin{theorem}[Generation by transpositions]
        The set of transpositions in $S_n$ generates $S_n$.
    \end{theorem}
    
    \begin{proof}
        We proceed by induction on $n$. In the base case, $n=2$, and the result is obvious. For the inductive step, suppose $S_{n-1}$ is generated by transpositions, and let $\sigma \in S_n$. If $\sigma(n) = n$, then $\sigma \in S_{n-1} \leq S_n$, and so we conclude by the inductive hypothesis. Otherwise, let $\tau = (n \; \sigma(n))$, so $\tau \sigma(n) = n$. Then, as before, it follows that $\tau \sigma \in S_{n-1}$, so $\tau \sigma$ is a product of transpositions -- multiplying by $\tau^{-1} = \tau$ on both sides, $\sigma$ is a product of transpositions as required.
    \end{proof}

    \lecture{Lecture 18 -- 19/11/25}
    
    \begin{lemma}[Adjacent transpositions]
        Any transposition $(i \; j)$ can be written as a product of an odd number of adjacent transpositions.
        \label{adjacent-transpositions}
    \end{lemma}
    
    \begin{proof}
        Suppose $j > i$, we proceed by induction on $j-i$. In the base case, $j=i+1$, and so $(i \; j) = (i \; i+1)$ as required. Otherwise, if $j>i+1$, we write
        \[
            (i\; j) = (j-1 \; j)(i \; j-1)(j-1 \; j)
        \]
        and the result follows by the inductive hypothesis applied to $(i \; j-1)$.
    \end{proof}
    
    So it turns out that $S_n$ is generated by adjacent transpositions. But the really important part of this is the notion of parity.

    \begin{lemma}
        If $\tau_1, \dots, \tau_k$ are all transpositions, and
        \[
            \sigma = \tau_1 \cdots \tau_k = e
        \]
        then $k$ is even.
    \end{lemma}

    \begin{definition}
        An unordered pair $\{i, j\} \subseteq \{1, \dots, n\}$ is called an \emph{inversion} of a permutation $\sigma$ if $i < j$ but $\sigma(i) > \sigma(j)$.
    \end{definition}
    
    
    
    \begin{proof}
        By \cref{adjacent-transpositions}, we may assume all $\tau_i$ are adjacent. We claim that, for any $\sigma = \tau_1 \cdots \tau_k$, $k \equiv \text{ the number of inversions of } \sigma \pmod 2$. We induct on $k$.
        
        In the base case, $k=0$, so $\sigma = e$ and we're done.
        
        For the inductive step, let $\sigma = \tau \sigma'$, where $\tau = \tau_1 = (l \; l+1)$ and $\sigma' = \tau_2 \cdots \tau_k$. Which pairs $\{i, j\}$ could be inversions of $\sigma'$ but not of $\sigma$, or vice versa? The only such pair would have to satisfy
        \begin{align*}
            \sigma'(i) = l 
            \sigma'(j) = l+1
        \end{align*}
        since $\tau$ fixes everything other than $l$ and $l+1$. Thus
        \begin{align*}
            \sigma'(i) = l < l + 1 = \sigma'(j) \\
            \sigma(i) = l+1 > l = \sigma(j)
        \end{align*}
        So, if $i < j$, then $\{i, j\}$ is an inversion for $\sigma$ but not for $\sigma'$, while if $i > j$, then $\{i, j\}$ is an inversion for $\sigma'$ but not for $\sigma$. In either case,
        \begin{align*}
            \#\text{inversions of }\sigma &= \#\text{inversions of }\sigma' \pm 1 \\
            &\equiv \#\text{inversions of }\sigma' + 1 \pmod 2
        \end{align*}
        as required. Then, the lemma follows, since $e$ has 0 inversions.
    \end{proof}
    
    This enables us to define the \emph{sign homomorphism}.
    
    \begin{theorem}[Sign homomorphism]
        The map
        \begin{align*}
            \sign: S_n &\to C_2 = \{\pm 1\} \\
            \tau_1 \cdots \tau_k &\mapsto (-1)^k
        \end{align*}
        where $\tau_i$ are transpositions, is a well-defined homomorphism.
    \end{theorem}
    
    \begin{proof}
        \begin{itemize}
            \item \emph{Well-defined}. Suppose we have
            \begin{align*}
                \tau_1 \cdots \tau_k &= \tau_1' \cdots \tau_l' \\
                \implies \tau_k \cdots \tau_1 \tau_1' \cdots \tau_l' &= e
            \end{align*}
            since all the transpositions are self-inverse. Hence, by the lemma, $k+l$ is even, so $k \equiv l \pmod 2$, and so $(-1)^k = (-1)^l$ as required.
            \item \emph{Homomorphism}. Clearly
            \begin{align*}
                \sign(\tau_1 \cdots \tau_k \tau_1' \cdots \tau_l') &= (-1)^{k+l} \\
                &= (-1)^k(-1)^l \\
                &= \sign(\tau_1 \cdots \tau_k)\sign(\tau_1' \cdots \tau_l')
            \end{align*}
        \end{itemize}
    \end{proof}
    
    \begin{definition}[Parity and the alternating group]
        If $\sign(\sigma) = 1$, then $\sigma$ is called \emph{even}, \emph{odd} otherwise. The subgroup
        \[
            A_n = \ker(\sign) \triangleleft S_n
        \]
        is called the \emph{alternating group} on $n$ elements -- it is the group of even permutations.
    \end{definition}
    
    \begin{example}
        Let's find $A_3$. We know
        \[
            S_3 = \{e, (1\; 2), (2 \; 3), (1 \; 3), (1 \; 2 \; 3), (3 \; 2 \; 1)\}
        \]
        and we have $(1 \; 2 \; 3) = (1 \; 2)(2 \; 3)$, and $(3 \; 2 \; 1) = (2 \; 3)(1 \; 2)$, so these are even. Thus,
        \[
            A_3 = \{e, (1 \; 2 \; 3), (3 \; 2 \; 1)\} \cong C_3
        \]
    \end{example}
    
    \begin{remark*}
        The cycle type makes it easy to determine a permutation's sign. Indeed,
        \[
            (a_1\; \cdots \; a_k) = (a_1 \; a_k) (a_1 \; a_{k-1}) \cdots (a_1 \; a_3)(a_1 \; a_2)
        \]
        so $(a_1 \; \cdots \; a_k)$ is a product of $k-1$ transpositions. Thus, $(a_1 \; \cdots \; a_k)$ is even if and only if $k$ is odd. More generally, a $(k_1, \dots, k_l)$-cycle is even if and only if $|\{k_i : k_i \text{ even}\}|$ is even. For example, $(1 \; 2)(3 \; 4)$ is even, but $(1 \; 2)(3 \; 4)(5 \; 6)$ is odd.
    \end{remark*}
    
    Let's apply this to study conjugacy in $S_n$ and $A_n$. $S_n$ is surprisingly easy!

    \begin{theorem}[Conjugacy in $S_n$]
        Two permutations $\sigma_1$, $\sigma_2 \in S_n$ are conjugate if and only if they have the same cycle type.
        \label{conjugacy-sn}
    \end{theorem}
    
    \begin{proof}
        \begin{itemize}
            \item[$\implies$] Suppose $\sigma_1 = (a_1^1 \; \cdots \; a_{l_1}^1)\cdots (a_1^k \; \cdots a_{l_k}^k)$ is a product of disjoint cycles. Thus $\sigma_1(a_j^i) = a_{j+1}^i$, where the addition in the subscript is modulo $l_i$. Since the cycle types are the same, $\sigma_2 = (b_1^1 \; \cdots \; b_{l_1}^1)\cdots (b_1^k \; \cdots b_{l_k}^k)$. Now consider the permutation $\tau$ given by $\tau(a_j^i) = b_j^i$. We have
            \begin{align*}
                \tau \sigma_1 \tau^{-1}(b_j^i) &= \tau\sigma_1(a_j^i) \\
                &= \tau(a^i_{j+1}) \\
                &= b^i_{j+1} \\
                &= \sigma_2(b^i_j)
            \end{align*}
            Therefore, $\sigma_2 = \tau\sigma_1\tau^{-1}$, so they are conjugate as required.
        \end{itemize}
        \lecture{Lecture 19 -- 21/11/25}
        \begin{itemize}
            \item[$\impliedby$] Suppose $\sigma_2 = \tau\sigma_1\tau^{-1}$. The above argument shows that, if
            \[
                \sigma_1 = (a_1^1 \; \cdots \; a_{l_1}^1)\cdots (a_1^k \; \cdots a_{l_k}^k)
            \]
            then
            \[
                \sigma_2 = (b_1^1 \; \cdots \; b_{l_1}^1)\cdots (b_1^k \; \cdots b_{l_k}^k)
            \]
            where $b_j^i = \tau(a_j^i)$ for all relevant $i$, $j$. Therefore, $\sigma_1$ and $\sigma_2$ have the same cycle type.
        \end{itemize}
    \end{proof}
    
    This makes it easy to count conjugacy classes in $S_n$.

    \begin{example}[$n=3$]
        Recall
        \[
            S_3 = \{e, (1\; 2), (2 \; 3), (1 \; 3), (1 \; 2 \; 3), (3 \; 2 \; 1)\}
        \]
        \begin{itemize}
            \item $|\ccl_{S_3}(1 \; 2)| = \binom{3}{2} = 3$;
            \item $|\ccl_{S_3}(1 \; 2 \; 3)| = 2 \cdot 1 = 2$.
        \end{itemize}
    \end{example}
    Let's also recall that \cref{orbit-stabiliser-theorem} (the orbit-stabiliser theorem) implies that
    \begin{align*}
        |C_{S_n}(\gamma)| = \frac{|S_n|}{|\ccl_{S_n}(\gamma)|} = \frac{n!}{|\ccl_{S_n}(\gamma)|}
    \end{align*}
    so we can also find the size of the centraliser pretty easily!

    \begin{example}[$n=4$]
        Now we want to do a similar calculation on the less familiar $S_4$.
        \begin{itemize}
            \item $|\ccl_{S_4}[(1 \; 2)(3 \; 4)]| = \frac{1}{2} \binom{4}{2} = 3$;
            \item $|C_{S_4}[(1 \; 2)(3 \; 4)]| = 24/3 = 8$. The cool thing about this is that we can start listing elements of the centraliser, and we know that, once we have 8 things, we're done. In particular, we have
            \begin{align*}
                C_{S_4}[(1 \; 2)(3 \; 4)] = \{&e, (1 \; 2)(3 \; 4), (1 \; 2), (3 \; 4), (1 \; 3)(2 \; 4), \\ &(1 \; 4)(2 \; 3), (1 \; 4 \; 2 \; 3), (1 \; 3 \; 2 \; 4)\}
            \end{align*}
        \end{itemize}
        Let's do this in full generality.
        \begin{center}
            \begin{tabular}{c | c | c}
                Typical element $\gamma$ & $|\ccl_{S_4}(\gamma)|$ & $|C_{S_4}(\gamma)|$ \\
                \hline
                $e$ & 1 & 24 \\
                $(1 \; 2)$ & $\binom 4 2 = 6$ & 4 \\
                $(1 \; 2)(3 \; 4)$ & $\frac 1 2 \binom 4 2 = 3$ & 8 \\
                $(1 \; 2 \; 3)$ & $2 \cdot \binom 4 3 = 8$ & 3 \\
                $(1 \; 2 \; 3 \; 4)$ & $3! = 6$ & 4
            \end{tabular}
        \end{center}
    \end{example}
    
    Counting conjugacy classes in $A_n$ can also be done with a little thought.

    \begin{lemma}[Conjugacy classes in $A_n$]
        Let $\gamma \in A_n \triangleleft S_n$.
        \begin{itemize}
            \item If an odd element of $S_n$ commutes with $\gamma$, then $\ccl_{A_n}(\gamma) = \ccl_{S_n}(\gamma)$.
            \item If every element of $S_n$ that commutes with $\gamma$ is even, then $\ccl_{S_n}(\gamma)$ splits into two:
            \[
                \ccl_{S_n}(\gamma) = \ccl_{A_n}(\gamma) \cup \ccl_{A_n}(\tau \gamma \tau^{-1})
            \]
            where $\tau$ is any transposition.
        \end{itemize}        
    \end{lemma}
    
    \begin{proof}
        By \cref{orbit-stabiliser-theorem} (orbit-stabiliser),
        \begin{align*}
            |S_n| &= |\ccl_{S_n}(\gamma)||C_{S_n}(\gamma)| \\
            |A_n| &= |\ccl_{A_n}(\gamma)||C_{A_n}(\gamma)|
        \end{align*}
        Then, since $|S_n| = 2|A_n|$, we have
        \begin{align*}
            |\ccl_{S_n}(\gamma)| = 2 \frac{|C_{A_n}(\gamma)|}{|C_{S_n}(\gamma)|} \cdot |\ccl_{A_n}(\gamma)| && (*)
        \end{align*}
        Now
        \[
            C_{A_n}(\gamma) = \ker(\left.\sign\right|_{C_{S_n}(\gamma)}: C_{S_n}(\gamma) \to C_2)
        \]
        The image has size 1 or 2, so, by \cref{isomorphism-theorem} (the isomorphism theorem),
        \[
            |C_{S_n(\gamma)} : C_{A_n(\gamma)}| = 1 \text{ or } 2
        \]
        \begin{itemize}
            \item If there is an odd element of $S_n$ that commmutes with $\gamma$, we hit $-1 \in C_2$, so the image has size 2 and hence
            \[
                |C_{S_n(\gamma)} : C_{A_n(\gamma)}| = 2
            \]
            Then, Lagrange and $(*)$ gives $|\ccl_{S_n}(\gamma)| = |\ccl_{A_n}(\gamma)|$ as required.
            \item If no odd elements of $S_n$ commute with $\gamma$, $|C_{S_n(\gamma)} : C_{A_n(\gamma)}| = 1$, $C_{S_n}(\gamma) = C_{A_n}(\gamma)$, so $(*)$ gives $|\ccl_{S_n}(\gamma)| = 2|\ccl_{A_n}(\gamma)|$, so $\ccl_{A_n}(\gamma)$ is half as big. Now take a transposition $\tau \in S_n$. Observe $\tau \gamma \tau^{-1} \in \ccl_{S_n}(\gamma)$. If $\tau \gamma \tau^{-1} \in \ccl_{A_n}(\gamma)$, then $\tau \gamma \tau^{-1} = \alpha \gamma \alpha^{-1}$ for some $\alpha \in A_n$. But then $(\tau \alpha) \gamma (\tau \alpha)^{-1} = \gamma$, while $\tau \alpha$ is odd, being the product of an odd $\tau$ and an even $\alpha$, so $C_{S_n}(\gamma)$ contains an odd element. But this is a contradiction. Hence $\tau \gamma \tau^{-1} \notin \ccl_{A_n}(\gamma)$. From this it's fairly clear that 
            \[
                \ccl_{S_n}(\gamma) = \ccl_{A_n}(\gamma) \cup \ccl_{A_n}(\tau \gamma \tau^{-1})
            \]
            as required.
        \end{itemize}
    \end{proof}
    
    \begin{example}[$A_4$]
        The even elements of $S_4$ are $e$, the $(2, 2)$-cycles and the 3-cycles.
        \begin{itemize}
            \item Since there are an odd number of $(2, 2)$-cycles, the conjugacy class of them cannot possibly split, so it remains intact.
            \item For a 3-cycle $\sigma = (1 \; 2 \; 3)$, $|C_{S_4}(1 \; 2 \; 3)| = 3$, so
            \[
                C_{S_4}(1 \; 2 \; 3) = \gen{(1 \; 2 \; 3)} \leq A_4
            \]
            so every element that commutes with $C_{S_4}(1 \; 2 \; 3)$ is even, and hence it splits.
        \end{itemize}
        Let us summarise:
        \begin{center}
            \begin{tabular}{c | c}
                Typical element $\gamma$ & $|\ccl_{A_4}(\gamma)|$ \\
                \hline
                $e$ & 1 \\
                $(1 \; 2)(3 \; 4)$ & 3 \\
                $(1 \; 2 \; 3)$ & 4 \\
                $(3 \; 2 \; 1)$ & 4 \\
            \end{tabular}
        \end{center}
    \end{example}

    \lecture{Lecture 20 -- 24/11/25}
    
    Finally, let's look at conjugacy in $S_5$ and $A_5$.

    \begin{example}[Conjugacy in $S_5$]
        We take the same approach as $S_4$.
        \begin{center}
            \begin{tabular}{c | c | c | c}
                Typical element $\gamma$ & $|\ccl_{S_5}(\gamma)|$ & $|C_{S_5}(\gamma)|$ & $\sign$ \\
                \hline
                $e$ & 1 & 120 & 1 \\
                $(1 \; 2)$ & $\binom 5 2 = 10$ & 12 & -1 \\
                $(1 \; 2 \; 3)$ & $2 \cdot \binom 5 3 = 20$ & 6 & 1 \\
                $(1 \; 2)(3 \; 4)$ & $\frac 1 2 \binom 5 2 \binom 3 2 = 15$ & 8 & 1 \\
                $(1 \; 2 \; 3)(4 \; 5)$ & $2 \binom 5 3 = 20$ & 6 & -1 \\
                $(1 \; 2 \; 3 \; 4)$ & $3! \cdot \binom 5 4 = 30$ & 4 & -1 \\
                $(1 \; 2 \; 3 \; 4 \; 5)$ & $4! = 24$ & 5 & 1 \\
            \end{tabular}
        \end{center}
    \end{example}
    
    \begin{example}[Conjugacy in $A_5$]
        \begin{itemize}
            \item The transposition $(4 \; 5)$ commutes with $(1 \; 2 \; 3)$ but is odd, so the conjugacy class doesn't split -- $\ccl_{A_5}(1 \; 2 \; 3) = \ccl_{S_5}(1 \; 2 \; 3)$.
            \item There are an odd number of double transpositions, so they definitely can't split -- $\ccl_{A_5}[(1 \; 2)(3 \; 4)] = \ccl_{S_5}[(1 \; 2)(3 \; 4)]$
            \item For 5-cycles, the centraliser is only of size 5 and, since it's of order 5, the elements of the centraliser must be exactly the cyclic subgroup it generates -- $C_{S_5}(1 \; 2 \; 3 \; 4 \; 5) = \gen {(1 \; 2 \; 3 \; 4 \; 5)} \leq A_5$. Hence, only even permutations commute with $(1 \; 2 \; 3 \; 4 \; 5)$, and so the conjugacy class splits.
        \end{itemize} \vspace{10pt}
        In summary,
        \begin{center}
            \begin{tabular}{c | c}
                Typical element $\gamma$ & $|\ccl_{A_5}(\gamma)|$ \\
                \hline
                $e$ & 1 \\
                $(1 \; 2 \; 3)$ & 20 \\
                $(1 \; 2)(3 \; 4)$ & 15 \\
                $(1 \; 2 \; 3 \; 4 \; 5)$ & 12 \\
                $(2 \; 1 \; 3 \; 4 \; 5)$ & 12
            \end{tabular}
        \end{center}
    \end{example}
    
    \begin{theorem}[$A_5$ is simple]
        $A_5$ is simple.
    \end{theorem}
    
    \begin{proof}
        Suppose $N \triangleleft A_5$. Then $N$ is a union of conjugacy classes in $A_5$. Hence, it is necessary to express a factor of 60 (by Lagrange) as a sum of the sizes of the conjugacy classes, and importantly we must have $\{e\}$ as one of our conjugacy classes. One can see that this cannot be done unless $N = \{e\}$ or $N = A_5$.
    \end{proof}
    
    \section{Matrix groups}

    Let $M_n(\RR) = \{n \text{ by } n \text{ matrices with entries in } \RR\}$. Matrix multiplication is associative, and the identity matrix provides an identity. However, not all matrices are invertible.
    
    \begin{lemma}
        $A \in M_n(\RR)$ has an inverse if and only if $\det A \neq 0$.
    \end{lemma}
    
    Now we actually can make our matrices into a group.

    \begin{definition}[General linear group]
        Let $\GL_n(\RR) = \{A \in M_n(\RR) : \det A \neq 0\}$. This is a group under matrix multiplication by the above, namely the \emph{general linear group}.
    \end{definition}
    
    \begin{lemma}
        For $A$, $B \in M_n(\RR)$, $\det(AB) = \det A \det B$.
    \end{lemma}
    
    \begin{corollary}
        It follows that $\det$ is a homomorphism:
        \begin{align*}
            \det: \GL_n(\RR) &\to \RR_\times = (\RR \setminus \{0\}, 1, \times) \\
            A &\mapsto \det A
        \end{align*}
    \end{corollary}
    
    \begin{definition}[Special linear group]
        The \emph{special linear group} $\SL_n(\RR) = \ker (\det) = \{A \in M_n(\RR): \det A = 1\}$.
    \end{definition}
    
    By \cref{isomorphism-theorem} (the isomorphism theorem), $\SL_n(\RR) \triangleleft \GL_n(\RR)$, and $\GL_n(\RR) / \SL_n(\RR) \cong \im(\det) \cong \RR_\times$ -- it's fairly clear to see that any real number can arise as a determinant.

    Furthermore, everything above still works with $\CC$, so we can talk about $\GL_n(\CC)$, $\SL_n(\CC)$, and analogously we have $\GL_n(\CC) / \SL_n(\CC) \cong \CC_\times$.
    
    \subsection{Change of basis}

    There is a natural action by conjugation of $\GL_n(\RR)$ on $M_n(\RR)$:
    \begin{align*}
        \GL_n(\RR) &\acts M_n(\RR) \\
        P(A) &:= PAP^{-1}
    \end{align*}
    
    \begin{proposition}
        Let $V$ be an $n$-dimensional vector space over $\RR$, and $\alpha: V \to V$ a linear map. If $A \in M_n(\RR)$ represents $\alpha$ in some basis, then the orbit
        \[
            \GL_n(\RR)A = \{PAP^{-1}: P \in \GL_n(\RR)\}
        \]
        consists of all the matrices representing $\alpha$ in any basis.
    \end{proposition}
    
    \begin{proof}
        A basis $\{\vec v_1, \dots, \vec v_n\}$ for $V$ defines an isomorphism of vector spaces
        \begin{align*}
            \phi: \RR^n &\to V \\
            (\lambda_1, \dots, \lambda_n) &\mapsto \sum_{i=1}^n \lambda_i \vec v_i
        \end{align*}
        The claim that $A$ represents $\alpha$ in this basis means that the following diagram commutes
        \[\begin{tikzcd}
            {\RR^n} & V \\
            {\RR^n} & V
            \arrow["\phi", from=1-1, to=1-2]
            \arrow["A"', from=1-1, to=2-1]
            \arrow["\alpha", from=1-2, to=2-2]
            \arrow["\phi"', from=2-1, to=2-2]
        \end{tikzcd}\]
        -- that is, $\alpha = \phi A \phi^{-1}$.
        
        \lecture{Lecture 21 -- 24/11/25}

        Likewise, another basis $\{\vec u_1, \dots, \vec u_n\}$ for $V$ corresponds to another isomorphism $\psi: \RR^n \to V$, and another matrix $B$ representing $\alpha$ in these coordinates, which again satisfies $\alpha = \psi B \psi^{-1}$. Therefore,
        \begin{align*}
            B &= \psi^{-1}\alpha \psi \\
            &= (\psi^{-1} \circ \phi) A(\phi^{-1} \circ \psi) \\
            &= PAP^{-1}
        \end{align*}
        where $P \in \GL_n(\RR)$ represents the linear map $\psi^{-1} \circ \phi: \RR^n \to \RR^n$ in the standard basis (note we know $P$ is invertible since $P^{-1}$ is the matrix representing $\phi^{-1} \circ \psi$). Thus, the set of all matrices representing $\alpha$ is contained in the orbit of $A$ under conjugation.

        Conversely, if $B = PAP^{-1}$ for some $P \in \GL_n(\RR)$, putting $\psi = \phi \circ P^{-1}: \RR^n \to V$ gives us a basis $\{\vec u_i = \psi(\vec e_i)\}$ for $V$. In this basis, $B$ represents $\alpha$.
    \end{proof}

    \subsection{M\"obius transformations: a new perspective}

    Recall that multiplication in $\mathcal M$ looks similar to multiplication of 2 by 2 matrices.
    
    \begin{proposition}
        Represent $\CC_\times$ as a matrix group by 
        \[
            \CC_\times = \left\{\begin{pmatrix}
             \lambda & 0 \\
             0 & \lambda \\
            \end{pmatrix} \in \GL_2(\CC): \lambda \in \CC_\times\right\}
        \]
        Then $\CC_\times \triangleleft \GL_2(\CC)$, and $\mathcal M \cong \GL_2(\CC)/\CC_\times$
    \end{proposition}
    
    \begin{proof}
        Let
        \begin{align*}
            \Phi: \GL_2(\CC) &\to \mathcal M \\
            \begin{pmatrix}
             a & b \\
             c & d \\
            \end{pmatrix} &\mapsto \left(z \mapsto \frac{az+b}{cz+d}\right)
        \end{align*}
        By our computation of multiplication in $\mathcal M$, $\Phi$ is a homomorphism, and it's clearly surjective, so $\im \Phi = \mathcal M$. Furthermore, a matrix $M$ is in $\ker \Phi$ if and only if its image fixes $0$, $1$ and $\infty$ Then
        \begin{itemize}
            \item $M$ fixes 0 $\implies b = 0$;
            \item $M$ fixes 1 $\implies a = d$;
            \item $M$ fixes $\infty$ $\implies c = 0$.
        \end{itemize}
        So $\ker \Phi = \CC_\times$. The result now follows by the isomorphism theorem.
    \end{proof}
    
    \begin{remark*}
        We get to simultaneously prove that $\CC_\times$ is normal and identify the quotient all just by looking at our homomorphism.
    \end{remark*}
    
    \subsection{Orthogonal groups}

    Let's write $||\cdot||$ for the usual notion of length in $\RR^n$ (the Euclidean norm) -- that is,
    \[
        ||\vec u|| = \sqrt{\sum_{i=1}^n u_i^2}
    \]

    \begin{definition}
        The \emph{$n$-dimensional orthogonal group}, $O(n)$, is the subgroup of $\GL_n(\RR)$ that preserves the Euclidean norm.
    \end{definition}
    In fact, the dot product is often more convenient to work with:
    \[
        \vec u \cdot \vec v = \sum_{i=1}^n u_iv_i
    \]
    and it turns out that matrices preserve all dot products if and only if they preserve dot products of vectors with themselves, i.e. the Euclidean norm.

    \begin{lemma}[Polarisation identity]
        For $\vec u$, $\vec v \in \RR^n$,
        \[
            2(\vec u \cdot \vec v) = ||\vec u||^2 + ||\vec v||^2 - ||\vec u - \vec v||^2
        \]
    \end{lemma}
    
    \begin{proof}
        Observe
        \begin{align*}
            ||\vec u - \vec v||^2 &= (\vec u - \vec v) \cdot (\vec u - \vec v) \\
            &= ||\vec u||^2 - 2(\vec u \cdot \vec v) + ||\vec v||^2
        \end{align*}
        whence the identity follows by rearrangement.
    \end{proof}
    In other words, length-preserving transformations also preserve angles. Importantly for us, this allows us to characterise $O(n)$ using the dot product.
    \begin{lemma}[$O(n)$ and the dot product]
        $O(n) = \{A \in \GL_n(\RR) : (A \vec x) \cdot (A \vec y) = \vec x \cdot \vec y \forall \vec x, \vec y \in \RR^n\}$.
    \end{lemma}
    
    \begin{proof}
        If $(A \vec x) \cdot (A \vec y) = \vec x \cdot \vec y$ for all $\vec x$, $\vec y \in \RR^n$, then, for any $\vec v \in \RR^n$,
        \begin{align*}
            ||A \vec v||^2 &= (A \vec v) \cdot (A \vec v) \\
            &= \vec v \cdot \vec v = ||\vec v||^2
        \end{align*}
        We may take square roots on both sides since the norm is positive. Conversely, if $A \in O(n)$, then, for any $\vec x$, $\vec y \in \RR^n$,
        \begin{align*}
            2(A \vec x) \cdot (A \vec y) &= ||A \vec x||^2 + ||A \vec y||^2 - ||A \vec x - A \vec y||^2 \\
            &= ||\vec x||^2 + ||\vec y||^2 + ||\vec x - \vec y||^2 \\
            &= 2 \vec x \cdot \vec y
        \end{align*}
    \end{proof}
    This quickly leads to a nice characterisation of matrices in $O(n)$.

    \begin{lemma}[Matrices in $O(n)$]
        Let $A \in M_n(\RR)$. The following are equivalent:
        \begin{enumerate}[label=\protect\circled{\arabic*}]
            \item $A \in O(n)$;
            \item the columns of $A$ form an orthonormal basis;
            \item $A^\intercal A = I$.
        \end{enumerate}
    \end{lemma}
    
    \begin{proof}
        Let $A = (a_{ij})$.
        \begin{itemize}[leftmargin=60px]
            \item[($1 \implies 2$)] Let $\{\vec e_1, \dots, \vec e_n\}$ be the standard basis for $\RR^n$. The $i^\text{th}$ column of $A$ is the column vector $A\vec e_i$. But then
            \[
                (A \vec e_i) \cdot (A \vec e_j) = \vec e_i \cdot \vec e_j = \delta_{ij}
            \]
            so the columns form an orthonormal basis, as required.
            \item[($2 \implies 3$)] We know $(A\vec e_i) \cdot (A \vec e_j) = \delta_{ij}$. Then
            \begin{align*}
                (A^\intercal A)_ij &= \vec e_i^\intercal (A^\intercal A) \vec e_j \\
                &= (A \vec e_i)^\intercal (A \vec e_j) = \delta_{ij}
            \end{align*}
            so $A^\intercal A = I$ as required. \\
            \item[($3 \implies 1$)] Suppose $\vec u$, $\vec v \in \RR^n$. Then
            \begin{align*}
                (A \vec u) \cdot (A \vec v) &= (A \vec u)^\intercal (A \vec v) \\
                &= \vec u^\intercal A^\intercal A \vec v \\
                &= \vec u^\intercal \vec v \\
                &= \vec u \cdot \vec v
            \end{align*}
            so $A \in O(n)$ as required.            
        \end{itemize}
    \end{proof}
    Recall that $\det A^\intercal = \det A$. Therefore,
    \[
        1 = \det I = \det (A^\intercal A) = (\det A)^2
    \]
    and so the determinant of any orthogonal matrix is $\pm 1$.

    \lecture{Lecture 22 -- 26/11/25}

    \begin{definition}
        The \emph{special orthogonal group} $SO(n) = O(n) \cap SL_n(\RR)$ -- the set of orthogonal matrices with determinant 1. Observe also that
        \[
            SO(n) = \ker(\det: O(n) \to \{\pm 1\})
        \]
        so $|O(n):SO(n)| = 2$.
    \end{definition}
    Examples of elements of $O(n)$ which are not in $SO(n)$ are provided by reflections.

    \begin{definition}[Reflections]
        Any $\vec v \in \RR \setminus \{0\}$ defines an \emph{orthogonal plane}
        \[
            \vec v^\perp = P_{\vec v} = \{\vec x \in \RR^n: \vec x \cdot \vec v = 0\}
        \]
        The \emph{reflection} in $P_{\vec v}$ is defined as
        \[
            S_{\vec v}(\vec x) = \vec x - \frac{2(\vec x \cdot \vec v)}{||\vec v||^2}\vec v
        \]
        DIAGRAM: any diagram that makes this definition make sense.
    \end{definition}
    
    \begin{remark}
        \begin{itemize}
            \item We will sometimes write $S_P$ for the reflection in the plane $P$.
            \item We may replace $\vec v$ by $\frac{\vec v}{||\vec v||}$ and so we can clearly assume $||\vec v|| = 1$ -- then $S_{\vec v}(\vec x) = \vec x - 2(\vec x \cdot \vec v)\vec v$.
        \end{itemize}
    \end{remark}
    
    \begin{lemma}
        $S_{\hvec v}^2 = I$ for any $\hvec v$, and $S_{\hvec v} \in O(n)$.
    \end{lemma}
    
    \begin{proof}
        Note that $S_{\hvec v}$ is linear in $\vec x$, so $S_{\hvec v} \in M_n(\RR)$. Now, observe
        \begin{align*}
            S_{\hvec v}(\vec x) \cdot \vec v &= (\vec x \cdot \hvec v) - 2(\vec x \cdot \hvec v)(\hvec v \cdot \hvec v) \\
            &= -\vec x \cdot \vec v
        \end{align*}
        Then
        \begin{align*}
            S_{\hvec v}^2(\vec x) &= S_{\hvec v}(\vec x) - 2(S_{\hvec v}(\vec x) \cdot \hvec v)\hvec v \\
            &= \vec x - 2(\vec x \cdot \hvec v)\hvec v + 2(\vec x \cdot \hvec v)\hvec v \\
            &= \vec x
        \end{align*}
        so $S_{\hvec v}^2 = I$ as required. Hence $S_{\hvec v}$ is non-singular, so $S_{\hvec v} \in \GL_n(\RR)$. Finally, for $\vec x \in \RR^n$,
        \begin{align*}
            ||S_{\hvec v}(\vec x)||^2 &= (\vec x - 2(\vec x \cdot \hvec v)\hvec v)(\vec x - 2(\vec x \cdot \hvec v)\hvec v) \\
            &= ||\vec x||^2 - 4(\vec x \cdot \hvec v)^2 + 4(\vec x \cdot \hvec v)^2 \\
            &= ||\vec x||^2
        \end{align*}
        so $S_{\hvec v} \in O(n)$ as required.        
    \end{proof}
    
    \begin{remark}
        Take a basis $\{\vec v_1, \dots, \vec v_{n-1}\}$ for $P_{\hvec v}$. Then, in the basis $\{\vec v_1, \dots, \vec v_{n-1}, \hvec v\}$, $S_{\hvec v}$ has matrix
        \[
            \begin{pmatrix}
             1 & \cdots & 0 & 0 \\
             \vdots & \ddots & 0 & 0 \\
             0 & 0 & 1 & 0 \\
             0 & 0 & 0 & -1 \\
            \end{pmatrix}
        \]
        so $\det S_{\hvec v} = 1$ and hence $S_{\hvec v} \notin SO(n)$.
    \end{remark}
    
    \begin{theorem}[Generation by reflections]
        Every $A \in O(n)$ is a product of at most $n$ reflections.
    \end{theorem}
    
    \begin{proof}
        We proceed by induction on $n$.

        \emph{Base case}. $n=1$. $O(1) = \{\pm 1\} = \gen {-1}$.

        \emph{Inductive step}. Let $\{\vec e_1, \dots, \vec e_n\}$ be the standard basis for $\RR^n$. Let $\vec v = \vec e_n - A \vec e_n$.

        DIAGRAM: $\vec e_n$ and $A \vec e_n$ shown as position vectors from the origin. They are connected, and perpendicular bisector drawn -- this is the plane we reflect in.

        Then $S_{\vec v}A \vec e_n = \vec e_n$. Since $S_{\vec v}A$ is orthogonal, it follows that $S_{\vec v}A$ preserves the plane $P_{\vec e_n} = \RR^{n-1} \times \{0\}$. By induction, there exist $\vec v_1, \dots, \vec v_{n-1} \in \RR^{n-1}$ such that $S_{\vec v}A = S_{\vec v_1} \cdots S_{\vec v_{n-1}}$ on $\RR^{n-1}$. Since both sides also fix $\vec e_n$, it follows they are also equal on $\RR^n$. Hence,
        \[  
            A = S_{\vec v}S_{\vec v_1} \cdots S_{\vec v_{n-1}}
        \]
        as required.
    \end{proof}
    
    Let us consider orthogonal transformations in low dimensions -- especially $n=2$.

    \begin{definition}[Rotations]
        A \emph{rotation} is
        \begin{itemize}
            \item in 2 dimensions, an orthogonal matrix that only fixes the origin;
            \item in three dimensions, an orthogonal matrix that only fixes a line.
        \end{itemize}
    \end{definition}
    
    

    \begin{lemma}[Elements of $O(2)$]
        Let $A \in O(2)$. If $A \notin SO(2)$, then $A$ is a reflection; if $A \in SO(2)$, then $A$ is a rotation about the origin.
    \end{lemma}
    
    \begin{proof}
        Recall that $\det S_{\vec v} = -1$, so the product of $k$ reflections has determinant $(-1)^k$. By the theorem, we may take $k \leq 2$.
        \begin{itemize}
            \item If $A \notin SO(2)$, then $k$ must be odd, so $k$ must be 1, so $A$ is a reflection.
            \item If $A \in SO(2)$, then $k$ is even so, unless $A = I$, $A = S_{\vec u}S_{\vec v}$ for some $\vec u$, $\vec v$ not parallel. We claim that $A$ only fixes $\vec 0$. Suppose we have $\vec x \neq \vec 0$ such that
            \begin{align*}
                S_{\vec u}S_{\vec v}\vec x = \vec x \iff S_{\vec u} \vec x = S_{\vec v} \vec x
            \end{align*}
            But $\vec v$ is parallel to $\vec x - S_{\vec v} \vec x$, and $\vec u$ is parallel to $\vec x - S_{\vec u} \vec x$, so $\vec u$ and $\vec v$ are parallel, which is a contradiction. Hence $A$ only fixes 0 as required, so it is a rotation.            
        \end{itemize}
    \end{proof}
    
    \begin{remark*}
        We've defined rotations to be orthogonal transformations that only fix the origin. But we can show that this is equivalent to other definitions. For example, let $A \in SO(2)$, so the columns of $A$ form an orthonormal basis. For the vectors to be perpendicular, we require
        \[
            A = \begin{pmatrix}
             a & b \\
             -b & a \\
            \end{pmatrix}
        \]
        and then we further require $a^2 + b^2 = 1$. But then $a = \cos \theta$, $b = \sin \theta$ for some $\theta$, so
        \[
            A = \begin{pmatrix}
             \cos \theta & \sin \theta \\
             -\sin \theta & \cos \theta \\
            \end{pmatrix}
        \]
        as required.
    \end{remark*}

    That concludes our survey of $O(2)$ and $SO(2)$.
    
    \begin{lemma}[Elements of $SO(3)$]
        If $A \in SO(3)$, $A$ is a rotation.
    \end{lemma}
    
    \begin{proof}
        As above, $A$ is a product of at most $3$ reflections so, for $\det A = 1$, $A = I$ or $A = S_{\vec u}S_{\vec v}$ for non-parallel $\vec u$, $\vec v$. It follows that
        \[
            P_{\vec u} \cap P_{\vec v} = l
        \]
        a line, so $A$ fixes $l$. Also, $S_{\vec u}S_{\vec v}x = x \implies S_{\vec u}x = S_{\vec v}x$ and, arguing as in the previous lemma, we find that either $S_{\vec u} = S_{\vec v}$, in which case $A = I$, or $x \in S_{\vec u} \cap S_{\vec v}$. Hence $A$ fixes only the line $l$, so $A$ is a rotation.
    \end{proof}
    
    \lecture{Lecture 23 -- 01/12/25}

    \section{Platonic solids}

    Amazingly, while there are infinitely many regular 2-dimensional polygons, there are only 5 \emph{Platonic solids} in 3 dimensions.

    \begin{definition}
        A convex polyhedron $X \subseteq \RR^3$ is a \emph{Platonic solid} if
        \begin{itemize}
            \item every face is a regular $n$-gon for some constant $n$;
            \item $G = \Isom(X)$ acts transitively on the faces;
            \item if $x \in X$ is the midpoint of a face, then $\Stab_G(x) \cong D_{2n}$.
        \end{itemize}
    \end{definition}
    
    Let us consider the 5 solids.

    \begin{itemize}
        \item \emph{The tetrahedron}. 4 triangular faces, 4 vertices.
        \item \emph{The cube/hexahedron}. 6 square faces, 8 vertices.
        \item \emph{The octahedron}. 8 triangular faces, 6 vertices.
        \item \emph{The dodecahedron}. 12 pentagonal faces, 20 vertices.
        \item \emph{The icosahedron}. 20 triangular faces, 12 vertices.
    \end{itemize}
    
    \begin{definition}
        Two solids $X$ and $Y$ are \emph{dual} if $Y$ can be constructed from $X$ by putting vertices in the centres of each face, and joining vertices in adjacent faces.
    \end{definition}
    
    The cube and the octahedron are dual (DIAGRAM), as are the dodecahedron and icosahedron. The tetrahedron is dual to itself. Furthermore, if $X$ and $Y$ are dual, $\Isom(X) \cong \Isom(Y)$, so we only have three groups to identify!
    
    \begin{example}[Tetrahedron]
        Let $G = \Isom(\text{tetrahedron})$. We want to look for a group action. By definition, $G$ acts transitively on the (4) faces, and each stabiliser is isomorphic to $D_6$. Thus, orbit-stabiliser gives
        \[
            |G| = 4 \times 6 = 24
        \]
        Furthermore, the action of $G$ on the 4 vertices defines a homomorphism $\theta: G \to S_4$. Let us prove that $\theta$ is injective. Suppose $\theta(g) = e$. Then $g$ fixes four vertices, which are not coplanar, and so $g = \id$ by the four-point lemma, meaning $\theta$ is injective as required. Then, by comparing sizes, we see $G \cong S_4$.

        Let's also identify the group of rotational symmetries:
        \[
            G_0 = G \cap SO(3)
        \]
        assuming our tetrahedron is centred on the origin.
    \end{example}
    
    \begin{lemma}[Uniqueness of $A_n$]
        If $H \leq S_n$ and $|S_n:H|=2$, then $H=A_n$.
    \end{lemma}
    
    \begin{proof}
        Because $|S_n:H|=2$, $H \triangleleft S_n$ and $S_n/H \cong C_2$. We therefore have a surjective homomorphism $\theta: S_n \to C_2$ with $\ker \theta = H$.
    \end{proof}
    
    \lecture{Lecture 24 -- 03/12/25}

    Finally, (non-examinable), 

    \begin{example}[Dodecahedron/icosahedron]
        Let $G = \Isom(\text{dodecahedron})$ and $G_0$ the rotational subgroup of $G$ (of index 2). Since a face has orbit of size 12 and stabiliser of size 10, $|G| = 120$, $|G_0| = 60$.

        By drawing diagonals on the faces, we may inscribe 5 cubes into the dodecahedron. Since the cubes are built symmetrically from the geometry of the dodecahedron, $G_0$ acts on them, and so we obtain a homomorphism
        \[  
            \theta: G_0 \to S_5
        \]
        The rotation around the axis through an opposite pair of vertices gives rise to a 3-cycle. There are ten different pairs of opposite vertices, so we get ten pairs of 3-cycles which are one another's inverses -- that way, we get all 20 3-cycles in $S_5$.

        We claim that $A_5$ is generated by 3-cycles. Since the 3-cycles form a conjugacy class, we know the subgroup generated by 3-cycles is normal in $A_5$ but, since $A_5$ is simple, it must be $A_5$ itself. Hence, $A_5 \leq \im \theta$, but
        \[
            60 = |G_0| \geq |\im \theta| \geq |A_5| = 60
        \]
        so $\theta$ is an isomorphism -- $G_0 \cong A_5$.

        Finally, just as for the cube, $-e$ is a symmetry of the dodecahedron which commutes with everything, so, by the direct product theorem, $G \cong A_5 \times C_2 \cong S_5$.
    \end{example}
    
    
    
    
\end{document}
